\documentclass[10pt,a4paper]{article}
\usepackage[T1]{fontenc}
\usepackage[utf8]{inputenc}
\usepackage[margin=1.6cm,top=2cm,bottom=2cm]{geometry}
\usepackage{amsmath,amssymb}
\usepackage{booktabs,longtable,array}
\usepackage{xcolor}
\usepackage{enumitem}
\usepackage[colorlinks=true,linkcolor=blue!50!black,urlcolor=blue!50!black]{hyperref}
\setlength{\parindent}{0pt}
\setlength{\parskip}{5pt}
\setlength{\tabcolsep}{3pt}
\setlength{\LTcapwidth}{\textwidth}
\setlist{nosep,leftmargin=*}
\definecolor{cst}{RGB}{0,100,0}
\definecolor{cin}{RGB}{150,85,0}
\definecolor{cmi}{RGB}{170,0,0}
\definecolor{cmy}{RGB}{70,40,140}
% The four evidence labels. The words carry the meaning, so they survive black and white printing.
\newcommand{\ST}[1]{\textcolor{cst}{\textbf{[STATED #1]}}}
\newcommand{\IN}[1]{\textcolor{cin}{\textbf{[INFERRED}: #1\textbf{]}}}
\newcommand{\MI}[1]{\textcolor{cmi}{\textbf{[MISSING}. Proposal: #1\textbf{]}}}
\newcommand{\MY}[1]{\textcolor{cmy}{\textbf{[MINE}: #1\textbf{]}}}
\newcolumntype{P}[1]{>{\raggedright\arraybackslash}p{#1}}
\newcommand{\Spread}{\mathit{Spread}}
\title{The Unintended Consequences of Corporate Bond ETFs:\\ Technical Specification for a Rebuild}
\author{Specification of Dannhauser and Hoseinzade, \emph{Review of Financial Studies} 35 (2022), 51--90\\
DOI 10.1093/rfs/hhab031, published online 22 March 2021}
\date{Specification only. No code, data or backtest. Written 1 November 2026.}
\begin{document}
\maketitle
\tableofcontents
\newpage

\section*{Start here}
\addcontentsline{toc}{section}{Start here}

\textbf{The result in one paragraph.} The paper asks whether ETFs are a distinct source of fragility in corporate bonds. It compares two bonds of the same issuer that faced different ETF selling pressure in the 2013 Taper Tantrum. The more exposed bond widens more, by 8.5 to 11.3 bps across the four significant months in the LQD setting. The gap then fades by December 2013.

\textbf{The rebuild is feasible but not fully determined.} A quant can rebuild all the regressions from this document. Twelve choices still need a decision, and Section~\ref{sec:gaps} ranks them. The first is the most serious, because the paper gives two different rules for picking the control bond.

\textbf{You cannot run the core tests on the data here.} This user holds TRACE for 15 issuers from 2023 to 2025. The core tests need all US corporate bond TRACE from April 2013 to May 2014. Section~\ref{sec:need} lists what each table needs.

\textbf{How to read the labels.} Every specification item carries one of these labels.
\begin{itemize}
\item \ST{p.~N}: the paper says it, on that page or table.
\item \IN{reason}: the paper does not say it, but I derive it from something it does say.
\item \MI{default}: the paper does not say it. The default is my proposal, not the paper's.
\item \MY{comment}: my own suggestion or view, never the authors'.
\end{itemize}

\textbf{Page numbers.} Page numbers refer to the advance-access PDF, which runs 1 to 40. \IN{the issue runs 51 to 90 and has the same 40 pages, so issue page = PDF page + 50. I did not see the issue pagination.}

\textbf{What I read.}
\begin{itemize}
\item All 40 pages of the PDF, as text and as page images for every table and figure.
\item The OUP supplementary file, \texttt{hhab031\_supplementary\_data.docx}. It holds Appendices A and B only. Its numbers match the PDF appendix in a token-by-token comparison. It adds no new method.
\item The SSRN version (abstract 3559433). I did not read it. A search confirmed that it exists, at 62 pages, but the page refused the fetch with HTTP 403. It may hold detail that the journal version cut.
\end{itemize}
The paper mentions no separate internet appendix. The OUP article page, as far as I could see it, listed only that one file and no replication code. Pages: \url{https://academic.oup.com/rfs/article/35/1/51/6179821} and \url{https://papers.ssrn.com/sol3/papers.cfm?abstract_id=3559433}.

% ------------------------------------------------------------------
\section{Research question, hypotheses and identification}\label{sec:rq}

\subsection{Research question}
Are the design features of ETFs a unique source of fragility for corporate bonds? \ST{p.~2} The paper names three features: intraday exchange trading, in-kind creation and redemption, and a liquidity gap that widens under stress. \ST{p.~2}

\subsection{Hypotheses}
The paper states these in prose. The labels H1 to H6 are mine, for reference.
\begin{longtable}{@{}P{0.6cm}P{7.6cm}P{3.5cm}P{4.6cm}@{}}
\toprule
& Hypothesis & Test & Label\\
\midrule\endhead
H1 & ETFs attract investors with greater liquidity needs than mutual funds. & Equation~\eqref{eq:vol}, Table~2 & \ST{p.~2, p.~12}\\
H2 & Exchange trading enables positive feedback trading, so ETF flows respond more to lagged returns. & Equation~\eqref{eq:fps}, Table~3 & \ST{p.~2, p.~13--14}\\
H3 & ETFs hold smaller liquidity buffers than mutual funds. & Table~1, panel A & \ST{p.~3, p.~16}\\
H4 & ETF outflows reach bonds more directly. The trade-to-flow slope is near one for ETFs and below one for active funds. & Equation~\eqref{eq:lou}, Table~4 & \ST{p.~3, p.~17}\\
H5 & Under stress, authorized participants (APs, the dealers who exchange baskets of bonds for ETF shares) keep the largest and most liquid ETFs near NAV through redemptions. This turns ETF selling into bond selling. & Figure~2, Tables~5, 6, A1 & \ST{p.~3--4, p.~21}\\
H6 & Redemptions create price pressure followed by reversal, not price discovery. & Equation~\eqref{eq:main}, Tables~8--10 & \ST{p.~23, p.~31}\\
\bottomrule
\end{longtable}

\subsection{Identification logic}
The Taper Tantrum is the shock. \ST{p.~3, p.~19} The paper treats it as a fundamental shock that produced a correlated demand shock. All bonds face the fundamental shock. Only the ETF exposure differs across bonds of one issuer.

The design removes confounds in four steps. \ST{p.~4, p.~23--24}
\begin{enumerate}
\item \textbf{Issuer fixed effects} absorb any issuer-level change in credit risk during the episode.
\item \textbf{Restricting the sample to bonds in ETF portfolios} removes selection into the ETF universe.
\item \textbf{Measuring spreads against a maturity-matched Treasury} removes term structure differences.
\item \textbf{Issue-level controls and mutual fund pressure} address bond-level rate sensitivity and mutual fund selling.
\end{enumerate}

\begin{longtable}{@{}P{2.2cm}P{4.7cm}P{5.5cm}P{4.1cm}@{}}
\caption*{\textbf{The variation that identifies the effect in each setting}}\\
\toprule
Setting & Identifying variation & Assumption it rests on & What it addresses\\
\midrule\endhead
1. LQD threshold (Table~8) & Two bonds of one issuer on either side of the \$750 million index size rule. & The size rule is the only thing that keeps control bonds out of LQD and its redemptions. \ST{p.~4, p.~25} & Selection into the most liquid ETF.\\
2. Trade-based pressure (Table~9) & Cross-bond differences in actual ETF holding changes, from funds in the extreme flow quartiles. & Within an issuer, ETF managers do not choose the bonds to sell on bond-level fundamentals. Setting 3 exists because this may fail. \IN{p.~28 raises manager discretion as the concern that setting 3 answers} & Whether the bond was truly in the baskets. \ST{p.~4, p.~26}\\
3. Hypothetical arbitrage (Table~10) & Pre-episode portfolio weights multiplied by each fund's net creation and redemption dollars. & Funds scale all holdings in proportion. Pre-episode weights are unrelated to later bond-specific news. \ST{p.~11, p.~28 for proportional scaling} \IN{the second part follows from using pre-period weights to avoid discretion} & Manager discretion in choosing baskets. \ST{p.~4, p.~28}\\
\bottomrule
\end{longtable}

\textbf{Pressure against price discovery.} The test uses the path of coefficients. Price discovery would need negative and significant coefficients after September 2013, as exposed bonds return first. Or it would need control bonds to catch up. Pressure predicts positive coefficients that fade. \ST{p.~31} The paper rejects catch-up from panel B of Figure~3, which is a visual argument and not a formal test. \ST{p.~31}

\textbf{Limit the authors set themselves.} They do not rule out broad price discovery through ETFs. They do not compare ETF prices to NAV in this test, or ETF bonds to non-ETF bonds. \ST{p.~5, p.~31}

\textbf{Dates of the episode.} The paper uses several windows, so name the one you mean.
\begin{itemize}
\item Onset: 22 May 2013, the Bernanke testimony. \ST{p.~19}
\item End of turmoil: 18 September 2013, the Fed decision to keep the programme. \ST{p.~19}
\item LQD redemption statistic: 21 May to 18 September. \ST{p.~23}
\item Flow window: May, June, July and August 2013. \ST{p.~11}
\item Creation and redemption window: 22 May to 31 August. \ST{p.~11}
\item Rating screen window: 21 May to 31 August. \ST{p.~11}
\item Outcome months: August 2013 to May 2014. The start follows Cella, Ellul and Giannetti (2013), and two proxies use flows through August. \ST{p.~25, fn.~27}
\end{itemize}

% ------------------------------------------------------------------
\section{Data}\label{sec:data}

\subsection{Sources, fields, frequency and filters}
The overall sample runs from January 2010 to December 2018. \ST{p.~6} The event tests use 2013 and 2014 only.

\begin{longtable}{@{}P{2.3cm}P{5cm}P{3cm}P{6.3cm}@{}}
\caption*{\textbf{Sources as used by the paper}}\\
\toprule
Source & Fields & Frequency, period & Filters and labels\\
\midrule\endfirsthead
\toprule Source & Fields & Frequency, period & Filters and labels\\ \midrule\endhead
CRSP Survivor-Bias-Free US Mutual Fund Database & Fund identification, holdings, returns, total net assets (TNA), turnover, expense ratio, share class data, Lipper code, fund name. & Monthly returns and TNA. Quarterly holdings. 2010--2018. \ST{p.~6, p.~11} & Keep funds with average corporate bond weight above 20\% and a taxable fixed-income Lipper asset code. \ST{p.~6} Keep $0.5<\text{TNA}_t/\text{TNA}_{t-1}<2$. \ST{p.~8} Name rules appear in Section~\ref{sec:steps}. \ST{p.~7, fn.~12}\\
Morningstar Mutual Fund Database & Investment class weights, average credit quality, average effective duration, average coupon. & \MI{monthly where available, carried forward within the quarter} & Matched to CRSP on fund CUSIP. \ST{p.~6}\\
CRSP ETF daily & Daily market price (\texttt{prc}) and daily NAV (\texttt{dnav}). & Daily. January 2012 to December 2014 in Figure~2. \ST{p.~10, p.~22} & \MI{all ETFs in the fund sample with both fields present}\\
Bloomberg & Daily ETF shares outstanding (\texttt{shrout}). & Daily. \ST{p.~10} \MI{cover 1 January 2013 to 30 June 2014, the widest test window} & None stated.\\
TRACE & Bond trades. The paper does not say whether it uses Standard or Enhanced TRACE. \MI{Enhanced TRACE, because it shows true size above the dissemination caps} & Trade level. 2010--2018 for liquidity. May 2013 to May 2014 for spreads. \IN{Tables~4 and 8--10 need these spans} & Dick-Nielsen (2009) error filter. Remove trades below \$100,000 par. Median and reversal filters (see below). Remove bonds under 1 year to maturity. \ST{p.~8}\\
Mergent Fixed Income Securities Database (FISD) & Bond characteristics, ratings from three agencies, amount outstanding, issue events. & \MI{reference data as of each date, with the amount history table} & Exclusions appear in Section~\ref{sec:steps}. \ST{p.~10--11}\\
Treasury rates & Daily maturity-matched Treasury rate. & Daily. \ST{p.~8} & Vendor \MI{Federal Reserve H.15 constant maturity yields, linearly interpolated to the bond's remaining maturity}\\
iBoxx USD Liquid Investment Grade Index rules & \$750 million minimum, 3 years minimum to maturity, issuer rules. & Rules as of 2013. \ST{p.~25} & The paper cites a Markit link for the issuer rules and does not list them. \MI{obtain the 2013 rule book. Without it, apply only the size and maturity rules}\\
SIFMA and Federal Reserve Board & Market size, issuance, holdings by investor type, dealer inventory. & Annual, 1996--2018. \ST{p.~6--7} & Figure~1 only. It is context, not a test.\\
\bottomrule
\end{longtable}

\textbf{TRACE price filters.} The median filter drops a trade whose price differs by more than 20\% from the daily median or the 9-day median. The reversal filter drops a trade whose price moves more than 10\% and then reverses. \ST{p.~8, fn.~13} \MI{for the reversal, flag a trade when its absolute price change from the previous trade exceeds 10\% and the next trade moves back in the opposite direction by more than 10\%. Check this against Edwards, Harris and Piwowar (2007), which I did not read}

\subsection{Reported sample counts}
The paper gives no attrition table. It reports only the counts below. A step-by-step count of bonds through the filters is \MI{produce your own attrition table at each step in Section~\ref{sec:steps}, and treat it as a deliverable}
\begin{longtable}{@{}P{9cm}P{3cm}P{4.8cm}@{}}
\toprule
Item & Count & Source\\
\midrule\endhead
Corporate bonds held by the funds before the onset & 10,131 & \ST{p.~10}\\
Share of these held by ETFs, active funds, index funds & 67\%, 96\%, 48\% & \ST{p.~10}\\
Distinct issuers in ETF holdings & 1,491 & \ST{p.~10}\\
ETF-holding issuers not held by active funds & 68 & \ST{p.~10}\\
Distinct index funds, and those with corporate weight over 50\% & 36 and 3 & \ST{p.~3, fn.~4}\\
Funds in the cross-section, without and with controls & 1,128 and 873 & \ST{Table 2, p.~13} \IN{each observation is one fund}\\
Share classes in the cross-section, without and with controls & 2,540 and 2,145 & \ST{Table 2, p.~13} \IN{each observation is one share class}\\
LQD setting, bonds per monthly regression & 151 to 171 & \ST{Table 8, p.~27}\\
Trade-based setting, bonds per monthly regression & 999 to 1,247 & \ST{Table 9, p.~29}\\
Hypothetical setting, bonds per monthly regression & 1,258 to 1,491 & \ST{Table 10, p.~30}\\
Number of ETFs, overall and per tercile & not reported & \MI{report your own. Do not treat it as a check}\\
Number of issuers in each setting & not reported & \IN{about 75 to 85 in Table 8, if most issuers have one treated and one control bond}\\
\bottomrule
\end{longtable}

\subsection{Plausible equivalents}
All of this subsection is \MY{my suggestions, not the paper's. Field names below come from my memory of these vendors, so verify them}
\begin{longtable}{@{}P{4cm}P{13cm}@{}}
\toprule
Paper's source & Equivalent\\
\midrule\endhead
CRSP mutual fund data & Morningstar Direct or LSEG Lipper for flows and TNA. SEC N-CSR and N-Q filings on EDGAR for 2013 holdings. N-PORT starts only in 2019 and cannot serve this window.\\
Morningstar portfolio statistics & Lipper fund data, or sponsor fact sheets for the few largest ETFs.\\
CRSP ETF price and NAV & Bloomberg \texttt{PX\_LAST} and \texttt{FUND\_NET\_ASSET\_VAL}. Sponsor premium and discount history, such as iShares downloads.\\
Bloomberg shares outstanding & Sponsor daily history, or the shares field in an ETF data vendor file. Public ICI data are aggregates and do not suffice.\\
Enhanced TRACE & WRDS Enhanced TRACE. The WRDS Bond Returns product gives monthly cleaned data, but its window will not match the last-three-days rule.\\
Mergent FISD & Bloomberg reference data, or LSEG reference data.\\
Treasury rates & FRED H.15. Gürkaynak, Sack and Wright zero curves as a robustness check.\\
iBoxx index rules and constituents & IHS Markit (now S\&P Global) archive. LQD historical holdings from iShares as a check on eligibility.\\
\bottomrule
\end{longtable}

\subsection{What the study needs, and what this user lacks}\label{sec:need}
\textbf{This user cannot run the bond-level tests.} TRACE here covers 15 issuers from 2023 to 2025, so the Taper Tantrum window is absent. The tests need the whole bond cross-section in 2013. \MY{I assessed feasibility only from the description given to me, and did not query the warehouse}

\begin{longtable}{@{}P{3.3cm}P{9.2cm}P{4.5cm}@{}}
\caption*{\textbf{Minimum data for each output. All rows are my derivation.}}\\
\toprule
Output & Needs & TRACE needed?\\
\midrule\endhead
Tables 1A, 2, 3 & CRSP funds and Morningstar, December 2008 to December 2018. The 12-month volatility needs a run-up year, or the first values fall in 2011. & No\\
Table 5 & CRSP monthly TNA and return, December 2012 to December 2013. & No\\
Figure 2, Tables 6, A1 & CRSP daily ETF price and NAV, 2012 to 2014. Daily shares, 2013. ETF bid and ask. Holdings of March and September 2013. & No\\
Table 4, columns without controls & CRSP quarterly holdings and monthly TNA, 2010 to 2018. & No\\
Table 4, columns with controls & As above, plus FISD amount outstanding and all-market TRACE from 2009 for the lagged Amihud measure. & Yes, 2009--2018\\
Table 1B & Holdings at the last report before May 2013. FISD. TRACE around May 2013. & Yes, about 2013\\
Tables 7--10, B1--B4, Figure 3 & TRACE from April 2013 to May 2014 (July 2014 for Figure 3), with January 2013 for Figure 3 panel A. FISD with ratings and events. Daily Treasury rates. Fund holdings from about March to September 2013. Shares from May to August 2013. Index rules. & Yes, all bonds\\
\bottomrule
\end{longtable}

% ------------------------------------------------------------------
\section{Variables}\label{sec:vars}

Subscripts: $i$ bond, $j$ issuer, $f$ fund, $t$ month, $d$ day, $q$ quarter. Winsorising: the paper states it for the bond spread only. \ST{p.~8--10} For every other variable, the paper is silent. \MI{no winsorising, so that the winsorising step cannot explain any gap to the targets. Report a 1\% and 99\% version as a robustness check}

\subsection{Fund flow and fund characteristics}
\begin{longtable}{@{}P{2.7cm}P{7.2cm}P{7.1cm}@{}}
\toprule
Variable & Formula and units & Timing and label\\
\midrule\endfirsthead
\toprule Variable & Formula and units & Timing and label\\ \midrule\endhead
$\mathit{Flow}_{f,t}$ & $\dfrac{\text{TNA}_{f,t}-\text{TNA}_{f,t-1}(1+R_{f,t})}{\text{TNA}_{f,t-1}}$. Percent in tables. Eq.~(1). & Monthly. \ST{p.~8} Dollar flow is the numerator, in \$ millions. \ST{Table 5, p.~20} Fund-level returns are TNA-weighted across share classes. \ST{p.~8}\\
$\mathit{Flow}_{f,\text{Tantrum}}$ & Geometric average of the four monthly flows. \IN{monthly flows can be negative, so the only workable form is $[\prod_{m}(1+\mathit{Flow}_{f,m})]^{1/4}-1$} & May to August 2013. \ST{p.~11}\\
Flow volatility & Standard deviation of monthly flow over 12 months. Percent. & Rolling monthly in the panel. Average per fund in the cross-section. \ST{p.~12--13} Rolling window \MI{trailing 12 months with all 12 observations required}\\
Expense ratio, turnover & CRSP. Percent in Table 1. \ST{p.~9} & Lagged 1 month in panels. Averaged and demeaned in the cross-section. \ST{p.~12--13} Regression scale \MI{percent, as in Table 1}\\
$\log(\text{age})$ & Natural log of years since inception. \IN{the text says log of age in years, and the paper does not state the base} & Inception date \MI{CRSP first offer date of the oldest share class}\\
$\log(\text{assets})$ & Log of TNA. & As for expense ratio. \ST{p.~12}\\
$\log(\text{family assets})$ & Log of total fund family assets. \ST{p.~12} & Family definition \MI{sum of TNA across CRSP funds with one management company code}\\
Credit quality & Portfolio average, numeric. \IN{lower means safer, because p.~8 says index funds have higher quality and they show the lowest score} & Morningstar. \ST{p.~6} Coding scale \MI{use the same rating scale as for bonds, and check against the Table 1 quartiles}\\
Effective duration, coupon & Portfolio average. Years and percent. \ST{Table 1, p.~9} & Lagged 1 month in Table 3. \ST{p.~14}\\
\% corporate, cash, Treasuries & Portfolio weights. Percent. \ST{Table 1, p.~9} & Morningstar. The ``average corporate weight'' screen uses a period \MI{average over all fund months in 2010--2018}\\
$\mathit{ETF}_f$, $\mathit{Index}_f$ & Dummies from CRSP identifiers and name strings. \ST{p.~6--7, fn.~12} & Fixed by fund. Vanguard ETF share classes form separate funds. \ST{p.~6--8}\\
\bottomrule
\end{longtable}

\subsection{ETF arbitrage variables}
\begin{longtable}{@{}P{2.7cm}P{7.2cm}P{7.1cm}@{}}
\toprule
Variable & Formula and units & Timing and label\\
\midrule\endhead
$\mathit{Deviation}$, daily & $(\text{prc}_{f,t,d}-\text{dnav}_{f,t,d})/\text{dnav}_{f,t,d}$. Percent in tables. Eq.~(3). Positive means premium. & Daily. \ST{p.~10}\\
$C/R$ dollars, daily & $\text{dnav}_{f,t,d}\,(\text{shrout}_{f,t,d}-\text{shrout}_{f,t,d-1})$. Dollars. Eq.~(4). Positive means net creation. & Daily. \ST{p.~10}\\
$C/R$ dollars, Tantrum & Sum of daily values. Dollars. & 22 May to 31 August 2013. \ST{p.~11} A version with lagged NAV gives the same results. \ST{fn.~31, p.~28}\\
Redemption dollars & Monthly sum of daily $C/R$ on days when the shares change is negative. & \ST{Figure 2, p.~22}\\
Deviation terciles & Sort ETFs by average daily deviation in the Tantrum. Bottom: Discount. Middle: Low Deviation. Top: Premium. \ST{p.~21--22} & Window \MI{22 May to 31 August 2013}\\
Net $C/R$ quintiles & Sort ETFs each month by total $C/R$ dollars. \ST{Table A1, p.~33} & Monthly, May to August 2013. Average volatility is the volatility of daily deviation. Scale \MI{report in the same units as deviation and compare ratios, since the Table A1 scale is unclear}\\
ETF bid-ask & (ask $-$ bid) / midpoint, end of day, averaged. Percent. \ST{Table 6, p.~23} & January to April, and May to August 2013.\\
Volume, ETF turnover, size & Secondary volume, turnover, total assets. \ST{Table 6, p.~23} & The row says shares and the note says dollars. \IN{use shares, because the row values look like daily share counts for ETFs of this size}\\
\bottomrule
\end{longtable}

\subsection{Pressure measures}
\textbf{ETF pressure}, Eq.~(5), then multiplied by minus one. \ST{p.~11}
\begin{equation}\label{eq:pressure}
\mathit{ETFpressure}_{ij}=-1\times\left[\frac{\sum_{f}\max(0,\Delta H_{ijf}\mid F_f>P_{75})}{AO_{ij}}-\frac{\sum_{f}\max(0,-\Delta H_{ijf}\mid F_f<P_{25})}{AO_{ij}}\right]
\end{equation}
\begin{itemize}
\item $\Delta H_{ijf}$ is the change in holdings, from the report before May to the first report after August 2013. \ST{p.~11}
\item $F_f$ is $\mathit{Flow}_{f,\text{Tantrum}}$. $P_{25}$ and $P_{75}$ are its quartiles. \ST{p.~11}
\item Positive values mean net selling pressure. \IN{the raw term is buys minus sells, and the paper flips the sign ``to aid in interpretation''}
\item Units are percent of amount outstanding. \IN{p.~26 reads the slope per 1\%, and Table 7 means are 0.277 and 0.458}
\item The numerator and denominator are both par. \MI{use reported par from holdings. If only market value exists, divide by price at the report date}
\item Quartiles use ETFs only for the ETF measure. \IN{p.~11 computes the measure ``for ETFs'', then a ``similar measure for all mutual funds''}
\item $AO_{ij}$ is dated $t-1$ in the paper. \IN{the date hardly matters, because the paper drops bonds whose amount outstanding changes}
\item A bond with no trade by an extreme-quartile fund gets zero, or may leave the sample. The paper says the measure is ``defined for bonds in the ETF universe that are in baskets''. \ST{p.~26} \IN{Table 9 has fewer bonds than Table 10, so bonds with no ETF trade probably leave the sample} \MI{keep a bond in Table 9 only if at least one ETF changed its holding between the two reports. Test the alternative that keeps zeros}
\end{itemize}

\textbf{MF pressure.} The same measure for all mutual funds, active and index combined, also times minus one. \ST{p.~11, Table 8 note} \IN{quartiles are computed across all mutual funds pooled, since the paper describes one combined measure}

\textbf{ETF hypothetical arbitrage exposure}, Eq.~(6), then multiplied by minus one. \ST{p.~11--12}
\begin{equation}\label{eq:arb}
\mathit{ETFArb}_{ij}=-1\times\sum_{f=1}^{F}\frac{w_{ijf,\text{pre}}\times C/R\,\text{dollars}_{f,\text{Tantrum}}}{AO_{ij}}
\end{equation}
\begin{itemize}
\item $w_{ijf,\text{pre}}$ is the weight of the bond in fund $f$ at the holdings report before the episode. \ST{p.~11} \MI{market value divided by fund TNA, as a fraction}
\item Units are percent of amount outstanding. \IN{p.~28 reads the slope per 1\%. So multiply by 100 if the weight is a fraction}
\item The numerator is at market value and the denominator at par. \IN{this follows from the two formulas. The paper does not correct for it}
\item The paper reports no summary statistics for this measure. \MI{report mean, standard deviation and quartiles, because you cannot check the scale otherwise}
\item The paper says the scaling removes the mechanical link that Wardlaw (2020) describes. \ST{p.~12}
\end{itemize}

\textbf{LQD bond dummy.} One for treatment bonds, zero for control bonds. \ST{p.~25}

\subsection{Bond spreads and liquidity}
\begin{longtable}{@{}P{2.7cm}P{7.2cm}P{7.1cm}@{}}
\toprule
Variable & Formula and units & Timing and label\\
\midrule\endfirsthead
\toprule Variable & Formula and units & Timing and label\\ \midrule\endhead
$\Spread_{ij,t}$ & Volume-weighted yield minus the daily maturity-matched Treasury rate. Percentage points. \IN{a coefficient of 0.113 is read as 11.3 bps on p.~26} & All trades in the last three full trading days of month $t$. \ST{p.~8, p.~25} The table notes say ``business days'', and fn.~14 settles it: half days do not count. \ST{fn.~14, p.~8}\\
Order of operations & \MI{compute a spread for each trade as yield minus the Treasury rate of that day, then take the par-volume-weighted mean} & Yield source \MI{the yield in Enhanced TRACE if present. Otherwise, compute yield to maturity from price and FISD terms}\\
Winsorising & 1\% and 99\%. \ST{p.~8--10} \IN{applied to the spread level, because the paper says the ``spread values'' are winsorised} & Cross-section or pooled \MI{by month, over all bonds that pass the filters, before differencing}\\
$\Spread_{ij,\text{May}}$ & As above. & Three days before the testimony. \ST{p.~25} \IN{17, 20 and 21 May 2013, because 22 May was a Wednesday and 18--19 May a weekend}\\
Spread change & $\Spread_{ij,t}-\Spread_{ij,\text{May}}$. Percentage points. & $t=$ August 2013 to May 2014. \ST{p.~25}\\
$\text{Amihud}_{t}$ & Monthly median of the daily value $\frac{1}{N_d}\sum_{d}\frac{|r_d|}{v_d}$. Eq.~(2). $r_d$: return between consecutive transactions. $v_d$: dollar volume in millions. & \ST{p.~10} Which trade gives $v_d$ \MI{the later trade of each pair} Minimum trades \MI{two trades a day}\\
Amihud scale & Percent per \$ million in Table 1 (mean 0.73). Fraction per \$ million in regressions. \IN{Table 7 shows 0.011 and 0.009. A slope of about 11 on 0.73 would imply 8 percentage points, which is not plausible} & \MY{treat this as a mandatory first scale check}\\
Amihud in Table 4 & Quarterly lagged 6-month moving average, divided by its standard deviation. \ST{Table 4 note, p.~18} & Which standard deviation \MI{cross-sectional, within quarter}\\
Roll & $2\sqrt{-\mathrm{cov}(R_t,R_{t-1})}$, using trades in the past 21 trading days, with at least four observations. Discard a positive covariance. & \ST{p.~33} Regression scale \IN{fraction, for the same reason as Amihud}\\
IRC & $(P_{\max}-P_{\min})/P_{\max}$ within an imputed round trip, which is a set of same-size trades on the same day. & \ST{p.~33} Monthly aggregate \MI{monthly mean. Weight round trips equally}\\
Turnover (Table 7) & Used but not defined. Values are 0.049 and 0.041. & \MI{monthly traded par divided by amount outstanding, as a fraction}\\
\bottomrule
\end{longtable}

\subsection{Bond controls}
The controls in Equation~\eqref{eq:main} are set to their May levels. \ST{p.~25}
\begin{longtable}{@{}P{3.2cm}P{6.8cm}P{7cm}@{}}
\toprule
Variable & Definition & Label\\
\midrule\endhead
Time to maturity & Years to maturity. & \ST{p.~25} Date \MI{21 May 2013}\\
Coupon & Percent. & \ST{p.~25}\\
$\log(\text{size})$ & Log of amount outstanding. & \ST{p.~25} Base \MI{natural log, in dollars. Units affect only the intercept}\\
Option dummy & One if the bond has option features. & \ST{p.~25} FISD fields \MI{one if the bond is callable or putable}. Make-whole only calls are a known issue. \MI{count them as callable, then test without them}\\
Amihud & May level. & \ST{p.~25} Window \MI{1 to 21 May 2013, to stay before the onset. Test the full month as well}\\
May spread & The base spread. & \ST{p.~25}\\
Rating & Numeric. AAA = 1, AA+ = 2, AA = 3, AA$-$ = 4, and so on. & \ST{Table 1 note, p.~9} Combining agencies \MI{mean of the most recent ratings available, rounded only for display, following fn.~15's remark about average ratings}\\
Ownership (Table 1) & Percent of the bond held by each fund type. & \IN{the row and the panel B note imply this}\\
Own (Table 4) & Previous ownership of the bond by the fund. & \ST{Table 4 note} Scale \MI{percent of amount outstanding}\\
\bottomrule
\end{longtable}

\subsection{Trade variable for Table 4}
$\mathit{Trade}_{ijf,q}$ is the percent change in holdings of bond $i$ by fund $f$ between reports in $q-1$ and $q$. \ST{p.~16}

$\mathit{Flow}_{f,q}$ is the quarterly flow. \ST{p.~16} \MI{apply Equation (1) to TNA between the two report dates}

The slope is tested against one, so the two must share units. \IN{a test of $\beta=1$ is only meaningful if both are percent}

A new position has no defined percent change. \MI{drop new positions, since the paper's new-bond statistic is an unreported test in fn.~20. A full exit counts as $-100\%$}
% ------------------------------------------------------------------
\section{Sample construction, in order}\label{sec:steps}

The paper does not give an order. It lists the rules. \IN{the order below is mine, chosen so that each step has its inputs ready}

\subsection*{A. Funds}
\begin{enumerate}
\item Start from all CRSP funds, January 2010 to December 2018. \ST{p.~6}
\item Keep funds with a taxable fixed-income Lipper asset code and an average corporate bond weight above 20\%. \ST{p.~6}
\item Delete exchange-traded notes, leveraged and inverse funds. Their lower-case names contain any of: plus, enhanced, inverse, 2x, 3x, ultra, 1.5x, 2.5x, exchange traded note, exchanged-traded note. \ST{fn.~12, p.~7}
\item Label ETFs. Use CRSP identifiers, and names with ishares, spdr, holdrs, streetracks, exchange traded or exchange-traded. \ST{fn.~12, p.~7}
\item Label index mutual funds among the remainder. Use CRSP identifiers, and names with index, idx or indx. \ST{fn.~12, p.~7}
\item Split Vanguard funds into a separate ETF fund and index fund. \ST{p.~6--8}
\item Merge Morningstar portfolio data on fund CUSIP. \ST{p.~6}
\item Aggregate to portfolio level. Sum share class assets and value-weight the other characteristics. \ST{p.~8}
\item Compute monthly flow. Drop observations with TNA ratio outside (0.5, 2). \ST{p.~8} The paper says this removes ``funds''. \MI{drop the fund-month only. A permanent drop would remove the fund's entire history}
\item For share class regressions only, drop share classes and ETFs under \$10 million and retirement share classes. \ST{fn.~17, p.~12} \MI{identify retirement classes with the CRSP retirement flag, or an R in the class name}
\end{enumerate}

\subsection*{B. Bond trades and reference data}
\begin{enumerate}[resume]
\item Clean TRACE with the Dick-Nielsen (2009) filter. \ST{p.~8}
\item Remove trades under \$100,000 par. \ST{p.~8}
\item Apply the median and reversal filters. \ST{p.~8}
\item Remove bonds with under 1 year to maturity. \ST{p.~8}
\item Merge FISD on bond CUSIP and build numeric ratings. \ST{p.~10} \IN{the merge key is the standard one, and the paper does not name it}
\item Compute spreads, winsorise, then compute Amihud, Roll and IRC. \ST{p.~8--10, p.~33}
\end{enumerate}

\subsection*{C. Event sample}
\begin{enumerate}[resume]
\item Remove bonds with under 1 year and 4 months to maturity. \ST{p.~11} Date \MI{measure at 21 May 2013}
\item Remove bonds whose rating from any of the three agencies crosses investment grade between 21 May and 31 August 2013. \ST{p.~11}
\item Remove sinking fund, defaulted, convertible, foreign currency, floating-rate and Yankee bonds. \ST{p.~11}
\item Remove bonds with amount outstanding changes: partial call, conversion, repurchase, pay-in-kind, tender or exchange. \ST{p.~11, fn.~16} Window \MI{22 May 2013 to 31 May 2014, because the outcome runs to May 2014. Test a window that ends on 31 August 2013}
\item Keep only bonds in ETF portfolios at the last report before May 2013. \ST{p.~4, p.~24, p.~28}
\end{enumerate}

\subsection*{D. Exposures}
\begin{enumerate}[resume]
\item Compute Tantrum flow for each fund. Find quartiles within ETFs and within mutual funds. \ST{p.~11}
\item Compute $\Delta H$ across the two reports, then ETF pressure and MF pressure. Flip the sign. \ST{p.~11}
\item Sum daily $C/R$ dollars from 22 May to 31 August. Compute ETFArb. Flip the sign. \ST{p.~11--12}
\end{enumerate}

\subsection*{E. The three regression samples}
\begin{enumerate}[resume]
\item \textbf{Setting 1, LQD.} Require at least 4 years to maturity before the Taper Tantrum, for both groups. \ST{p.~25}
\item Treatment: within each issuer, take the smallest bond at or above \$750 million. If two tie, keep both. \ST{p.~25, fn.~29} The bond must be investment grade. \ST{p.~25} The paper does not say whether the bond must truly be in LQD at the report date. The text and the Table 8 note say ``meets the eligibility requirements''. \MI{use eligibility as the main rule, then a version that also requires an actual LQD holding}
\item Control: within the same issuer, a bond below \$750 million, held by another ETF. \ST{p.~25} The paper then gives two rules for picking it. See the first row of the ranked table in Section~\ref{sec:gaps}. \MI{closest May Amihud, one control per treatment, without replacement}
\item The alternative matches use Roll and IRC (Table B1), and time to maturity (Table B2). \ST{p.~26}
\item \textbf{Setting 2, trade-based.} All event-sample bonds with a defined ETF pressure. \ST{p.~26}
\item \textbf{Setting 3, hypothetical.} All event-sample bonds held by ETFs. \ST{p.~28}
\item For each month $t$, keep bonds with a May spread and a spread in month $t$. \IN{the bond count changes each month, so the sample is not balanced} Keep issuers with two or more bonds in that month. \IN{singletons add nothing under issuer fixed effects, and their presence does not affect the slope estimates, but it does affect $N$} \MI{report $N$ both ways}
\end{enumerate}

% ------------------------------------------------------------------
\section{Regressions and tests}\label{sec:reg}

Standard errors in all panels are clustered, and t-statistics appear in the tables. Stars: * $p<.1$, ** $p<.05$, *** $p<.01$. \ST{table notes}

\subsection{Flow volatility (Table 2)}
\begin{equation}\label{eq:vol}
\mathrm{SD}_{f,t}=\alpha+\beta_1\mathit{ETF}_f+\beta_2\mathit{Index}_f+\gamma' X_{f,t}+\varepsilon_{f,t}
\end{equation}
\begin{itemize}
\item The paper writes $\beta_2$ twice in Eq.~(7). I use $\gamma$ for the controls. \IN{a typographical slip, since the terms differ}
\item Controls: expense ratio, turnover, $\log$ age, $\log$ TNA, $\log$ family assets. \ST{p.~12}
\item Cross-section (columns 1, 2, 5, 6): dependent variable is the fund's average 12-month volatility. Controls are averages over the sample, then demeaned. No clustering or fixed effects are reported. \ST{p.~12--13} \MI{use heteroskedasticity-robust standard errors}
\item Panel (columns 3, 4, 7, 8): monthly volatility, controls lagged 1 month, month fixed effects, clustered by fund or share class. \ST{p.~13}
\item Columns 1--4 use funds. Columns 5--8 use share classes. \ST{p.~13}
\item Columns 1, 3, 5, 7 have only the ETF dummy. \ST{Table 2}
\end{itemize}

\subsection{Flow-performance sensitivity (Table 3)}
\begin{equation}\label{eq:fps}
\mathit{Flow}_{f,t}=\lambda_t+\delta(\mathit{Ret}_{f,t-1}\times\mathit{ETF}_f)+\beta_1\mathit{Ret}_{f,t-1}+\beta_2\mathit{ETF}_f+\beta_3' X_{f,t-1}+\varepsilon_{f,t}
\end{equation}
\begin{itemize}
\item Coefficient of interest: $\delta$. \ST{p.~14}
\item Month fixed effects. Clusters by fund or share class. \ST{p.~14--15}
\item Columns 1 and 5: no controls. Columns 2 and 6 add duration, credit quality, $\log$ size, $\log$ age, expense ratio, turnover. Columns 3 and 7 add three lags of flow. Columns 4 and 8 add 12-month flow volatility. \ST{Table 3}
\item The comparison group is all mutual funds, active and index. \IN{the note defines the dummy against ``mutual funds'', and no index dummy appears}
\end{itemize}

\subsection{Trade response to flows (Table 4)}
\begin{equation}\label{eq:lou}
\mathit{Trade}_{ijf,q}=\alpha+\tau_q+\beta_1\mathit{Flow}_{f,q}+\beta_2' X_{ijf,q}+\varepsilon_{ijf,q}
\end{equation}
\begin{itemize}
\item Run separately for ETFs, active funds and index funds, and separately for outflows and inflows. That gives 12 columns. \ST{p.~16--18}
\item Controls: lagged Amihud, ownership, and each interacted with flow. \ST{p.~16--17, Table 4}
\item Quarter fixed effects. Clusters by fund. \ST{Table 4 note}
\item Test: $H_0:\beta_1=1$. The table gives p-values. \ST{Table 4}
\item Outflow means $\mathit{Flow}<0$ and inflow means $\mathit{Flow}>0$. \IN{the words imply this. The paper does not state how it treats zero flow}
\item \MY{with the interaction terms, $\beta_1$ is the slope at Amihud = 0 and Own = 0. It is not the slope at the mean, because the paper scales Amihud without saying it is demeaned. Compare the two before reading 0.946 as the average response}
\end{itemize}

\subsection{Descriptive evidence on ETF arbitrage}
\begin{longtable}{@{}P{2.3cm}P{9.2cm}P{5.5cm}@{}}
\toprule
Output & Construction & Label\\
\midrule\endfirsthead
\toprule Output & Construction & Label\\ \midrule\endhead
Figure 1 & Market size and issuance. Holdings by investor type and dealer inventory. Annual. & \ST{p.~6--7}. Context only.\\
Table 1 & Panel A: fund characteristics, 2010--2018. Panel B: bond holdings at the report before onset, by fund type. & \ST{p.~9}\\
Table 5 & Monthly distribution of flows in 2013 by fund type. Percent and \$ millions. & \ST{p.~20}\\
Figure 2 & Panel A: mean deviation by tercile, monthly, January 2012 to December 2014. Panel B: total assets by tercile, May to August 2013. Panel C: total redemption dollars by tercile, same months. & \ST{p.~22}. I read the axis of panel A, and it covers the caption's dates, not those on p.~21.\\
Table 6 & Discount against Low Deviation terciles. January to April, and May to August 2013. Means and p-values of differences. & \ST{p.~23} Test \MI{two-sample t-test with unequal variances}\\
Table A1 & Monthly quintiles on net $C/R$, May to August 2013. & \ST{p.~33}\\
Table 7 & Treatment against control bond means, with t-values and p-values. & \ST{p.~26} Test \MI{two-sample t-test}. Date \MI{May levels, as for the regression controls}\\
Figure 3 & Panel A: mean spread difference of LQD held against same-issuer not held. Panel B: mean spreads. January 2013 to July 2014. & \ST{p.~24}. Holdings date \MI{last report before May 2013}. Weights \MI{equal-weighted bonds}. Sample \IN{the caption says bonds held by LQD, which is wider than the matched sample in Table 8}\\
\bottomrule
\end{longtable}

\subsection{The main event regressions (Tables 8, 9, 10)}
All three settings are one equation, Eq.~(10). \ST{p.~25} Each runs as a separate cross-section per month, for ten months. \ST{p.~24--25, table headers}
\begin{equation}\label{eq:main}
\bigl(\Spread_t-\Spread_{\text{May}}\bigr)_{ij}=\lambda_j+\beta_1\mathit{ETFproxy}_{ij}+\beta_2\mathit{MFpressure}_{ij}+\beta_3' X_{ij,\text{May}}+\varepsilon_{ij,t}
\end{equation}
for $t\in\{\text{Aug 2013},\dots,\text{May 2014}\}$. The notes to Tables 9 and 10 add $\alpha$, which is only a normalisation alongside $\lambda_j$. \IN{an intercept is collinear with a full set of issuer effects}

\begin{longtable}{@{}P{2.2cm}P{4.3cm}P{5.2cm}P{5.3cm}@{}}
\caption*{\textbf{The three versions of Equation~\eqref{eq:main}}}\\
\toprule
& Setting 1: Table 8 & Setting 2: Table 9 & Setting 3: Table 10\\
\midrule\endhead
$\mathit{ETFproxy}$ & LQD bond dummy & ETF pressure, Eq.~\eqref{eq:pressure} & ETFArb, Eq.~\eqref{eq:arb}\\
MF pressure & Included & Included & Excluded \ST{p.~28}\\
Controls $X$ & Time to maturity, coupon, Amihud, May spread. No option dummy because it has no variation. No size because it correlates with treatment. \ST{p.~26} & Time to maturity, coupon, $\log$ size, option dummy, Amihud, May spread. \ST{Table 9} & Same as setting 2. \ST{Table 10}\\
Fixed effects & Issuer & Issuer & Issuer\\
Clustering & Issuer \ST{Table 8} & Issuer \ST{Table 9} & Issuer \ST{Table 10}\\
Bond sample & Matched treatment and control & Bonds with a defined ETF pressure & All ETF-held event bonds\\
Robustness & Table B1 (match on Roll, IRC), Table B2 (match on maturity, with each proxy) & Table B3 (Roll, IRC as the proxy in $X$) & Table B4 (Roll, IRC as the proxy in $X$)\\
\bottomrule
\end{longtable}

\textbf{Economic scaling.} The paper reports 20.35 bps as the peak price effect of a 1\% pressure rise, ``accounting for average duration''. \ST{p.~26--28} \IN{$0.045\times 100\times 4.52=20.34$ bps, using the peak Table 9 coefficient and the ETF mean effective duration in Table 1. The fourth digit may come from an unrounded input}

\textbf{Subsamples.} The paper splits funds by type and flow sign in Table 4, and bonds by setting. It splits the main regressions by no other dimension. It does not split them by rating, maturity or liquidity. \IN{no such table appears}

\textbf{No pre-period test.} The paper runs no placebo regression before the event. \IN{no such table appears} See Section~\ref{sec:threats}.

% ------------------------------------------------------------------
\section{Target numbers}\label{sec:targets}

All numbers below come from the PDF. A script parsed most tables from the PDF text and confirmed that every number occurs there. I also compared the rendered tables with the page images. The tables keep the paper's numbers, so Table 8 here is Table 8 there. T-statistics are in parentheses. The paper reports no standard errors. \MY{derive them as $|\text{coefficient}/t|$. For example, the Table 8 August standard error is $0.113/2.67=0.042$}

{\scriptsize\setlength{\tabcolsep}{4pt}
\begin{longtable}{@{}lrrrrrr@{}}
\caption*{\textbf{Table 1, p.~9. Means and medians.} The paper also gives SD and quartiles.}\\
\toprule
 & \multicolumn{2}{c}{ETFs} & \multicolumn{2}{c}{Active funds} & \multicolumn{2}{c}{Index funds}\\
 & Mean & 50\% & Mean & 50\% & Mean & 50\%\\
\midrule\endhead
\multicolumn{7}{@{}l}{\emph{A. Funds}}\\
Flow monthly (\%) & $1.92$ & $0.13$ & $0.00$ & $-0.23$ & $0.36$ & $0.31$\\
Return monthly (\%) & $0.20$ & $0.16$ & $0.19$ & $0.19$ & $0.12$ & $0.09$\\
Total net assets (\$ bln) & $3.11$ & $0.36$ & $2.08$ & $0.47$ & $10.65$ & $1.70$\\
Expense ratio (\%) & $0.35$ & $0.30$ & $0.90$ & $0.91$ & $0.32$ & $0.26$\\
Turnover (\%) & $45.70$ & $25.00$ & $114.82$ & $70.00$ & $102.42$ & $77.00$\\
Age (years) & $4.72$ & $4.26$ & $8.46$ & $6.00$ & $7.72$ & $5.90$\\
Credit quality & $8.76$ & $9.00$ & $10.23$ & $9.00$ & $4.26$ & $3.00$\\
Effective duration (years) & $4.52$ & $4.39$ & $4.07$ & $4.24$ & $5.78$ & $5.53$\\
Coupon (\%) & $4.48$ & $4.31$ & $4.51$ & $4.18$ & $3.38$ & $3.28$\\
\% in corporate bonds & $73.88$ & $87.62$ & $56.09$ & $51.04$ & $37.99$ & $26.38$\\
\% in cash & $2.83$ & $0.91$ & $4.88$ & $3.77$ & $2.77$ & $1.33$\\
\% in Treasuries & $12.48$ & $2.13$ & $17.17$ & $11.24$ & $38.42$ & $42.06$\\
12m flow vol, fund & $7.13$ & $5.14$ & $3.75$ & $2.26$ & $3.69$ & $2.27$\\
12m flow vol, share class & & & $4.88$ & $2.70$ & $4.38$ & $2.56$\\
\midrule\multicolumn{7}{@{}l}{\emph{B. Bonds held, report before onset}}\\
Spread (\%) & $1.66$ & $1.25$ & $2.24$ & $1.39$ & $1.26$ & $1.12$\\
Size (\$ mln) & $728.30$ & $500.00$ & $637.67$ & $499.98$ & $734.32$ & $500.00$\\
Coupon (\%) & $5.33$ & $5.50$ & $5.66$ & $5.75$ & $4.91$ & $5.14$\\
Time to maturity (years) & $9.62$ & $6.20$ & $9.00$ & $5.91$ & $10.54$ & $6.82$\\
Rating & $8.88$ & $8.00$ & $9.75$ & $9.00$ & $7.64$ & $7.00$\\
Amihud & $0.73$ & $0.43$ & $0.85$ & $0.47$ & $0.78$ & $0.47$\\
Roll & $0.45$ & $0.34$ & $1.04$ & $0.36$ & $0.49$ & $0.37$\\
Imputed round-trip cost & $0.21$ & $0.17$ & $0.25$ & $0.18$ & $0.21$ & $0.17$\\
Ownership (\%) & $2.21$ & $1.99$ & $9.47$ & $6.14$ & $2.00$ & $1.82$\\
\bottomrule
\end{longtable}}
{\scriptsize\setlength{\tabcolsep}{4pt}
\begin{longtable}{@{}l*{8}{r}@{}}
\caption*{\textbf{Table 2, p.~13. Flow volatility. Columns 1--4 funds, 5--8 share classes. Columns 1, 2, 5, 6 cross-section.}}\\
\toprule
 & (1) & (2) & (3) & (4) & (5) & (6) & (7) & (8)\\
\midrule\endhead
ETF & $3.312^{***}$ & $1.454^{***}$ & $3.361^{***}$ & $2.155^{***}$ & $2.048^{***}$ & $0.908^{*}$ & $1.374^{***}$ & $0.871^{*}$\\
 & $(10.19)$ & $(2.59)$ & $(9.36)$ & $(4.60)$ & $(5.14)$ & $(1.73)$ & $(3.74)$ & $(1.91)$\\
Index &  & $0.983^{*}$ &  & $0.523$ &  & $-0.344$ &  & $-0.089$\\
 &  & $(1.92)$ &  & $(1.51)$ &  & $(-0.84)$ &  & $(-0.26)$\\
Expense ratio &  & $-0.014$ &  & $-0.005$ &  & $-0.004$ &  & $0.003$\\
 &  & $(-0.78)$ &  & $(-0.27)$ &  & $(-0.16)$ &  & $(0.20)$\\
Turnover ratio &  & $-0.017$ &  & $-0.003$ &  & $0.025$ &  & $0.004$\\
 &  & $(-1.01)$ &  & $(-0.41)$ &  & $(1.36)$ &  & $(0.76)$\\
log(age) &  & $-0.842^{***}$ &  & $-0.411^{***}$ &  & $-1.879^{***}$ &  & $-1.768^{***}$\\
 &  & $(-6.59)$ &  & $(-6.96)$ &  & $(-17.44)$ &  & $(-19.22)$\\
log(assets) &  & $-0.571^{***}$ &  & $-0.686^{***}$ &  & $-0.266^{***}$ &  & $-0.459^{***}$\\
 &  & $(-9.36)$ &  & $(-12.48)$ &  & $(-6.03)$ &  & $(-15.51)$\\
log(family assets) &  & $0.149^{***}$ &  & $0.182^{***}$ &  & $0.007$ &  & $0.032$\\
 &  & $(3.04)$ &  & $(3.59)$ &  & $(0.20)$ &  & $(0.95)$\\
Date fixed effects & No & No & Yes & Yes & No & No & Yes & Yes\\
Observations & $1{,}128$ & $873$ & $102{,}266$ & $73{,}376$ & $2{,}540$ & $2{,}145$ & $267{,}410$ & $230{,}584$\\
Adj. $R^2$ & $.084$ & $.234$ & $.049$ & $.117$ & $.009$ & $.173$ & $.013$ & $.121$\\
\bottomrule
\end{longtable}}
{\scriptsize\setlength{\tabcolsep}{4pt}
\begin{longtable}{@{}l*{8}{r}@{}}
\caption*{\textbf{Table 3, p.~15. Flow-performance sensitivity. Columns 1--4 funds, 5--8 share classes. Date fixed effects in all.}}\\
\toprule
 & (1) & (2) & (3) & (4) & (5) & (6) & (7) & (8)\\
\midrule\endhead
Ret(t-1) x ETF & $0.844^{***}$ & $0.565^{***}$ & $0.475^{***}$ & $0.484^{***}$ & $0.769^{***}$ & $0.570^{***}$ & $0.516^{***}$ & $0.524^{***}$\\
 & $(6.21)$ & $(3.45)$ & $(2.71)$ & $(2.76)$ & $(5.86)$ & $(4.05)$ & $(3.37)$ & $(3.43)$\\
Ret(t-1) & $0.265^{***}$ & $0.335^{***}$ & $0.240^{***}$ & $0.241^{***}$ & $0.347^{***}$ & $0.401^{***}$ & $0.267^{***}$ & $0.266^{***}$\\
 & $(5.71)$ & $(7.31)$ & $(5.66)$ & $(5.65)$ & $(11.26)$ & $(12.22)$ & $(8.65)$ & $(8.51)$\\
ETF & $2.467^{***}$ & $1.885^{***}$ & $0.986^{***}$ & $0.575^{***}$ & $2.004^{***}$ & $0.324$ & $0.259^{*}$ & $0.048$\\
 & $(10.25)$ & $(7.82)$ & $(6.42)$ & $(3.69)$ & $(9.88)$ & $(1.50)$ & $(1.81)$ & $(0.31)$\\
Effective duration &  & $-0.096^{***}$ & $-0.053^{***}$ & $-0.042^{***}$ &  & $-0.010$ & $-0.003$ & $-0.001$\\
 &  & $(-3.82)$ & $(-3.78)$ & $(-2.93)$ &  & $(-0.61)$ & $(-0.29)$ & $(-0.06)$\\
Credit quality &  & $-0.022^{*}$ & $-0.016^{**}$ & $-0.022^{***}$ &  & $-0.025^{***}$ & $-0.022^{***}$ & $-0.032^{***}$\\
 &  & $(-1.68)$ & $(-2.12)$ & $(-2.82)$ &  & $(-2.70)$ & $(-3.89)$ & $(-5.49)$\\
log(size) &  & $0.078^{***}$ & $0.041^{**}$ & $0.109^{***}$ &  & $0.205^{***}$ & $0.111^{***}$ & $0.160^{***}$\\
 &  & $(2.87)$ & $(2.47)$ & $(6.60)$ &  & $(10.74)$ & $(9.13)$ & $(13.16)$\\
log(age) &  & $-0.303^{***}$ & $-0.154^{***}$ & $-0.119^{***}$ &  & $-1.516^{***}$ & $-0.777^{***}$ & $-0.591^{***}$\\
 &  & $(-7.50)$ & $(-7.00)$ & $(-5.22)$ &  & $(-23.31)$ & $(-18.66)$ & $(-14.73)$\\
Expense ratio &  & $0.003$ & $0.001$ & $-0.001$ &  & $0.009^{**}$ & $0.004^{*}$ & $0.002$\\
 &  & $(0.89)$ & $(0.49)$ & $(-0.26)$ &  & $(2.57)$ & $(1.81)$ & $(0.84)$\\
Turnover &  & $-0.009$ & $-0.004$ & $-0.005$ &  & $-0.002$ & $0.001$ & $-0.000$\\
 &  & $(-1.19)$ & $(-1.04)$ & $(-1.11)$ &  & $(-0.49)$ & $(0.60)$ & $(-0.20)$\\
Flow(t-1) &  &  & $0.169^{***}$ & $0.159^{***}$ &  &  & $0.168^{***}$ & $0.155^{***}$\\
 &  &  & $(13.79)$ & $(13.02)$ &  &  & $(23.43)$ & $(21.40)$\\
Flow(t-2) &  &  & $0.136^{***}$ & $0.126^{***}$ &  &  & $0.111^{***}$ & $0.097^{***}$\\
 &  &  & $(11.09)$ & $(10.23)$ &  &  & $(16.27)$ & $(14.17)$\\
Flow(t-3) &  &  & $0.118^{***}$ & $0.106^{***}$ &  &  & $0.108^{***}$ & $0.092^{***}$\\
 &  &  & $(14.02)$ & $(12.59)$ &  &  & $(18.88)$ & $(15.94)$\\
Vol flow &  &  &  & $0.137^{***}$ &  &  &  & $0.156^{***}$\\
 &  &  &  & $(9.81)$ &  &  &  & $(18.09)$\\
Observations & $97{,}822$ & $58{,}828$ & $58{,}045$ & $58{,}045$ & $204{,}536$ & $140{,}257$ & $133{,}966$ & $133{,}966$\\
Adj. $R^2$ & $.033$ & $.034$ & $.122$ & $.130$ & $.017$ & $.042$ & $.110$ & $.121$\\
\bottomrule
\end{longtable}}
{\tiny\setlength{\tabcolsep}{2pt}
\begin{longtable}{@{}l*{12}{r}@{}}
\caption*{\textbf{Table 4, p.~18. Trade response to flows. Quarter fixed effects in all. Clustered by fund.}}\\
\toprule
 & \multicolumn{4}{c}{ETFs} & \multicolumn{4}{c}{Active funds} & \multicolumn{4}{c}{Index funds}\\
 & Out & Out & In & In & Out & Out & In & In & Out & Out & In & In\\
\midrule\endhead
Flow & $0.872^{***}$ & $0.946^{***}$ & $0.362^{***}$ & $0.411^{***}$ & $0.426^{***}$ & $0.440^{***}$ & $0.216^{***}$ & $0.267^{***}$ & $0.153$ & $0.156$ & $0.103^{*}$ & $0.113$\\
 & $(15.75)$ & $(15.13)$ & $(8.92)$ & $(8.24)$ & $(6.20)$ & $(4.94)$ & $(8.02)$ & $(8.90)$ & $(1.37)$ & $(1.29)$ & $(1.73)$ & $(1.54)$\\
Amihud &  & $-0.466^{**}$ &  & $-0.452^{*}$ &  & $0.663^{***}$ &  & $0.584^{***}$ &  & $0.002$ &  & $-0.307^{***}$\\
 &  & $(-2.16)$ &  & $(-1.82)$ &  & $(5.04)$ &  & $(4.94)$ &  & $(0.02)$ &  & $(-3.47)$\\
Own &  & $-3.159^{***}$ &  & $-4.289^{***}$ &  & $0.827^{***}$ &  & $0.685^{***}$ &  & $-3.978^{*}$ &  & $-0.312$\\
 &  & $(-3.15)$ &  & $(-5.94)$ &  & $(3.62)$ &  & $(4.47)$ &  & $(-1.82)$ &  & $(-0.83)$\\
Amihud x Flow &  & $-0.054^{***}$ &  & $-0.036^{**}$ &  & $-0.013$ &  & $-0.014$ &  & $0.014$ &  & $0.002$\\
 &  & $(-3.57)$ &  & $(-2.22)$ &  & $(-0.76)$ &  & $(-1.45)$ &  & $(0.95)$ &  & $(0.17)$\\
Own x Flow &  & $0.143$ &  & $0.108^{**}$ &  & $0.007$ &  & $-0.042^{***}$ &  & $-0.091$ &  & $-0.009$\\
 &  & $(1.36)$ &  & $(2.10)$ &  & $(0.19)$ &  & $(-4.67)$ &  & $(-0.21)$ &  & $(-0.15)$\\
Obs. & $179{,}191$ & $142{,}330$ & $346{,}560$ & $278{,}521$ & $1{,}628{,}939$ & $1{,}017{,}761$ & $1{,}618{,}023$ & $1{,}011{,}777$ & $217{,}522$ & $155{,}510$ & $423{,}827$ & $303{,}088$\\
Adj. $R^2$ & $.186$ & $.168$ & $.124$ & $.104$ & $.114$ & $.065$ & $.105$ & $.055$ & $.091$ & $.111$ & $.091$ & $.095$\\
p [$\beta=1$] & $.02$ & $.39$ & $.00$ & $.00$ & $.00$ & $.00$ & $.00$ & $.00$ & $.00$ & $.00$ & $.00$ & $.00$\\
\bottomrule
\end{longtable}}
{\scriptsize\setlength{\tabcolsep}{2pt}
\begin{longtable}{@{}l*{12}{r}@{}}
\caption*{\textbf{Table 5, p.~20. Monthly flows in 2013.  \$ rows are \$ millions.}}\\
\toprule
 & \multicolumn{4}{c}{ETFs} & \multicolumn{4}{c}{Active funds} & \multicolumn{4}{c}{Index funds}\\
 & Mean & p10 & p50 & p90 & Mean & p10 & p50 & p90 & Mean & p10 & p50 & p90\\
\midrule\endhead
Jan \% & $7.46$ & $-0.44$ & $2.17$ & $20.96$ & $1.83$ & $-2.02$ & $0.40$ & $7.11$ & $-0.63$ & $-6.21$ & $0.10$ & $5.20$\\
Jan \$ & $30.82$ & $-0.18$ & $3.33$ & $162.53$ & $20.99$ & $-13.76$ & $1.08$ & $65.32$ & $-20.55$ & $-64.21$ & $0.69$ & $48.25$\\
Feb \% & $6.79$ & $-2.33$ & $2.45$ & $21.04$ & $1.02$ & $-2.09$ & $0.10$ & $5.26$ & $1.05$ & $-2.90$ & $0.99$ & $5.02$\\
Feb \$ & $12.66$ & $-18.30$ & $5.00$ & $168.56$ & $11.39$ & $-16.92$ & $0.09$ & $47.38$ & $186.39$ & $-112.78$ & $3.61$ & $83.91$\\
Mar \% & $4.36$ & $-1.19$ & $1.47$ & $14.50$ & $1.31$ & $-2.22$ & $0.27$ & $5.72$ & $0.82$ & $-2.24$ & $0.54$ & $3.34$\\
Mar \$ & $40.34$ & $-13.05$ & $2.99$ & $136.25$ & $12.22$ & $-15.92$ & $0.37$ & $43.86$ & $56.36$ & $-51.79$ & $3.56$ & $58.10$\\
Apr \% & $2.81$ & $-1.18$ & $-0.02$ & $11.79$ & $1.19$ & $-2.11$ & $0.23$ & $6.34$ & $1.02$ & $-1.10$ & $0.40$ & $3.64$\\
Apr \$ & $41.05$ & $-15.93$ & $-0.01$ & $289.38$ & $9.04$ & $-15.57$ & $0.37$ & $42.69$ & $65.86$ & $-9.06$ & $4.57$ & $99.78$\\
May \% & $6.00$ & $-0.84$ & $2.72$ & $15.72$ & $1.50$ & $-2.67$ & $-0.04$ & $5.60$ & $0.66$ & $-2.08$ & $0.50$ & $5.40$\\
May \$ & $24.92$ & $-4.40$ & $5.02$ & $129.67$ & $10.12$ & $-22.16$ & $-0.02$ & $47.66$ & $-255.91$ & $-33.98$ & $1.76$ & $101.07$\\
Jun \% & $4.10$ & $-7.82$ & $-0.21$ & $17.49$ & $-1.05$ & $-6.59$ & $-1.08$ & $3.74$ & $0.08$ & $-2.64$ & $-0.30$ & $2.83$\\
Jun \$ & $-54.08$ & $-89.58$ & $-0.04$ & $34.15$ & $-41.85$ & $-119.81$ & $-2.32$ & $15.66$ & $-59.08$ & $-214.86$ & $-2.05$ & $81.74$\\
Jul \% & $3.07$ & $-7.66$ & $-0.10$ & $19.99$ & $0.68$ & $-4.04$ & $-0.37$ & $6.66$ & $-0.43$ & $-4.26$ & $0.02$ & $3.71$\\
Jul \$ & $30.01$ & $-38.94$ & $-0.02$ & $128.74$ & $-11.53$ & $-46.23$ & $-0.52$ & $46.48$ & $-8.16$ & $-192.72$ & $0.55$ & $75.13$\\
Aug \% & $2.56$ & $-6.37$ & $-0.15$ & $12.60$ & $0.18$ & $-4.76$ & $-0.82$ & $5.44$ & $0.94$ & $-2.44$ & $-0.09$ & $4.69$\\
Aug \$ & $-23.51$ & $-107.48$ & $-0.04$ & $56.56$ & $-13.67$ & $-59.10$ & $-1.58$ & $41.84$ & $81.30$ & $-30.93$ & $-0.70$ & $138.94$\\
Sep \% & $3.51$ & $-5.68$ & $0.26$ & $16.62$ & $0.39$ & $-3.54$ & $-0.13$ & $5.48$ & $-0.01$ & $-4.06$ & $-0.13$ & $3.57$\\
Sep \$ & $21.77$ & $-25.81$ & $0.06$ & $101.18$ & $-4.49$ & $-31.56$ & $-0.06$ & $46.16$ & $43.74$ & $-65.71$ & $-2.04$ & $83.28$\\
Oct \% & $3.49$ & $-3.31$ & $-0.11$ & $12.15$ & $0.31$ & $-3.69$ & $-0.20$ & $4.35$ & $1.02$ & $-3.98$ & $-0.28$ & $4.19$\\
Oct \$ & $17.57$ & $-35.83$ & $-0.02$ & $92.80$ & $-7.70$ & $-36.92$ & $-0.23$ & $36.07$ & $28.65$ & $-129.40$ & $-0.65$ & $55.38$\\
Nov \% & $2.36$ & $-3.78$ & $0.18$ & $12.12$ & $0.17$ & $-3.43$ & $-0.26$ & $3.84$ & $1.63$ & $-1.34$ & $0.29$ & $3.47$\\
Nov \$ & $-1.66$ & $-28.96$ & $0.02$ & $43.75$ & $-9.92$ & $-37.92$ & $-0.48$ & $28.17$ & $43.43$ & $-41.22$ & $2.25$ & $108.29$\\
Dec \% & $2.96$ & $-3.67$ & $0.20$ & $12.30$ & $0.22$ & $-3.80$ & $-0.35$ & $4.16$ & $-0.57$ & $-6.02$ & $0.20$ & $4.12$\\
Dec \$ & $-2.36$ & $-47.53$ & $0.03$ & $38.04$ & $-6.39$ & $-34.53$ & $-0.58$ & $33.16$ & $85.84$ & $-56.81$ & $1.87$ & $106.81$\\
\bottomrule
\end{longtable}}
{\scriptsize\setlength{\tabcolsep}{4pt}
\begin{longtable}{@{}l*{8}{r}@{}}
\caption*{\textbf{Table 6, p.~23. Discount and Low Deviation ETFs.}}\\
\toprule
 & \multicolumn{4}{c}{January--April 2013} & \multicolumn{4}{c}{May--August 2013}\\
 & Disc. & Low dev. & Diff. & p & Disc. & Low dev. & Diff. & p\\
\midrule\endhead
Size (\$ mln) & $904.89$ & $3{,}795.31$ & $-2{,}890.42$ & $<.0001$ & $908.31$ & $3{,}694.40$ & $-2{,}786.09$ & $<.0001$\\
Volume (shares) & $135{,}797$ & $587{,}130$ & $-451{,}332$ & $<.0001$ & $181{,}785$ & $673{,}265$ & $-491{,}480$ & $.002$\\
ETF turnover (\%) & $85.28$ & $40.17$ & $45.10$ & $.137$ & $58.92$ & $39.27$ & $19.65$ & $.351$\\
ETF bid-ask (\%) & $0.57$ & $0.16$ & $0.41$ & $.000$ & $0.77$ & $0.21$ & $0.56$ & $<.0001$\\
Average rating & $8.00$ & $7.33$ & $0.67$ & $.532$ & $8.20$ & $7.35$ & $0.85$ & $.382$\\
Average effective duration & $6.51$ & $5.15$ & $1.36$ & $.219$ & $6.44$ & $5.05$ & $1.39$ & $.164$\\
Average coupon & $4.61$ & $4.62$ & $-0.01$ & $.986$ & $4.48$ & $4.55$ & $-0.07$ & $.874$\\
\bottomrule
\end{longtable}}
{\scriptsize\setlength{\tabcolsep}{6pt}
\begin{longtable}{@{}lrrrrr@{}}
\caption*{\textbf{Table 7, p.~26. Treatment and control bonds. Signs are as printed.}}\\
\toprule
 & Non-LQD & LQD & Difference & t-value & p-value\\
\midrule\endhead
ETF pressure & $0.277$ & $0.458$ & $0.181$ & $2.50$ & $.013$\\
Amount outstanding (\$ mln) & $415.832$ & $1{,}122.307$ & $706.475$ & $15.30$ & $.000$\\
Rating & $7.415$ & $7.326$ & $-0.089$ & $-0.42$ & $.674$\\
Time to maturity & $15.273$ & $16.488$ & $1.215$ & $1.14$ & $.269$\\
Coupon & $5.224$ & $5.441$ & $0.216$ & $1.14$ & $.254$\\
MF pressure & $1.082$ & $0.757$ & $-0.326$ & $0.95$ & $.341$\\
Amihud & $0.011$ & $0.009$ & $-0.002$ & $1.64$ & $.101$\\
Turnover & $0.049$ & $0.041$ & $-0.008$ & $1.07$ & $.285$\\
\bottomrule
\end{longtable}}
\newpage
{\scriptsize\setlength{\tabcolsep}{2pt}
\begin{longtable}{@{}l*{10}{r}@{}}
\caption*{\textbf{Table 8, p.~27. LQD threshold. Issuer fixed effects, clustered by issuer.}}\\
\toprule
 & Aug 13 & Sep 13 & Oct 13 & Nov 13 & Dec 13 & Jan 14 & Feb 14 & Mar 14 & Apr 14 & May 14\\
\midrule\endhead
LQD bond & $0.113^{**}$ & $0.085^{**}$ & $0.112^{***}$ & $0.102^{**}$ & $0.052$ & $0.029$ & $-0.003$ & $0.006$ & $0.004$ & $0.035$\\
 & $(2.67)$ & $(2.20)$ & $(3.47)$ & $(2.51)$ & $(1.09)$ & $(0.80)$ & $(-0.08)$ & $(0.18)$ & $(0.09)$ & $(0.94)$\\
MF pressure & $-0.023$ & $-0.038$ & $-0.016$ & $0.000$ & $-0.006$ & $0.016$ & $-0.004$ & $0.006$ & $0.011$ & $0.001$\\
 & $(-0.64)$ & $(-1.47)$ & $(-0.67)$ & $(0.01)$ & $(-0.15)$ & $(0.57)$ & $(-0.17)$ & $(0.26)$ & $(0.35)$ & $(0.05)$\\
Time to maturity & $0.001$ & $0.000$ & $0.001$ & $-0.002$ & $0.000$ & $0.003$ & $0.004$ & $0.011^{***}$ & $0.015^{***}$ & $0.018^{***}$\\
 & $(0.23)$ & $(0.07)$ & $(0.46)$ & $(-0.48)$ & $(0.09)$ & $(0.92)$ & $(1.46)$ & $(4.20)$ & $(3.82)$ & $(5.05)$\\
Coupon & $0.004$ & $-0.025$ & $-0.004$ & $-0.014$ & $-0.002$ & $-0.001$ & $-0.014$ & $0.007$ & $0.002$ & $-0.016$\\
 & $(0.23)$ & $(-1.46)$ & $(-0.30)$ & $(-0.91)$ & $(-0.11)$ & $(-0.03)$ & $(-0.93)$ & $(0.53)$ & $(0.10)$ & $(-1.02)$\\
Amihud & $3.971$ & $-0.293$ & $-0.595$ & $6.049$ & $11.018^{***}$ & $11.642^{***}$ & $5.639^{**}$ & $9.049^{***}$ & $7.818^{**}$ & $7.772^{**}$\\
 & $(0.92)$ & $(-0.08)$ & $(-0.18)$ & $(1.58)$ & $(2.94)$ & $(3.58)$ & $(2.13)$ & $(2.85)$ & $(2.24)$ & $(2.60)$\\
May spread & $-0.388^{**}$ & $-0.176$ & $-0.280^{**}$ & $-0.229$ & $-0.432^{***}$ & $-0.324^{**}$ & $-0.239^{**}$ & $-0.374^{***}$ & $-0.415^{***}$ & $-0.478^{***}$\\
 & $(-2.41)$ & $(-1.32)$ & $(-2.42)$ & $(-1.63)$ & $(-2.73)$ & $(-2.57)$ & $(-2.10)$ & $(-4.06)$ & $(-3.07)$ & $(-3.08)$\\
Observations & 151 & 164 & 165 & 158 & 156 & 160 & 171 & 170 & 167 & 164\\
Adj.\ $R^2$ & .852 & .876 & .904 & .897 & .887 & .894 & .837 & .883 & .849 & .897\\
\bottomrule
\end{longtable}}
{\scriptsize\setlength{\tabcolsep}{2pt}
\begin{longtable}{@{}l*{10}{r}@{}}
\caption*{\textbf{Table 9, p.~29. Trade-based pressure. Issuer fixed effects, clustered by issuer.}}\\
\toprule
 & Aug 13 & Sep 13 & Oct 13 & Nov 13 & Dec 13 & Jan 14 & Feb 14 & Mar 14 & Apr 14 & May 14\\
\midrule\endhead
ETF pressure & $0.011$ & $0.039^{**}$ & $0.041^{***}$ & $0.045^{**}$ & $0.045^{**}$ & $0.038^{**}$ & $0.014$ & $0.020$ & $-0.007$ & $0.003$\\
 & $(0.77)$ & $(2.52)$ & $(2.67)$ & $(2.30)$ & $(2.25)$ & $(2.05)$ & $(0.92)$ & $(1.18)$ & $(-0.38)$ & $(0.20)$\\
MF pressure & $-0.018^{***}$ & $-0.012^{*}$ & $-0.003$ & $0.004$ & $-0.005$ & $-0.003$ & $0.010$ & $0.008$ & $0.008$ & $0.007$\\
 & $(-2.73)$ & $(-1.76)$ & $(-0.38)$ & $(0.50)$ & $(-0.65)$ & $(-0.32)$ & $(1.60)$ & $(1.15)$ & $(1.01)$ & $(0.96)$\\
Time to maturity & $-0.002^{*}$ & $-0.000$ & $-0.001$ & $0.000$ & $0.002$ & $0.006^{***}$ & $0.004^{***}$ & $0.008^{***}$ & $0.008^{***}$ & $0.011^{***}$\\
 & $(-1.73)$ & $(-0.01)$ & $(-1.01)$ & $(0.05)$ & $(0.80)$ & $(3.51)$ & $(3.08)$ & $(5.47)$ & $(4.89)$ & $(9.26)$\\
Coupon & $-0.025^{***}$ & $-0.026^{***}$ & $-0.021^{***}$ & $-0.034^{***}$ & $-0.030^{***}$ & $-0.036^{***}$ & $-0.033^{***}$ & $-0.034^{***}$ & $-0.037^{***}$ & $-0.040^{***}$\\
 & $(-3.64)$ & $(-3.73)$ & $(-3.12)$ & $(-3.98)$ & $(-3.41)$ & $(-4.59)$ & $(-5.25)$ & $(-4.70)$ & $(-4.54)$ & $(-6.77)$\\
log(size) & $0.086^{***}$ & $0.103^{***}$ & $0.069^{***}$ & $0.105^{***}$ & $0.071^{**}$ & $0.086^{***}$ & $0.046^{**}$ & $0.044^{*}$ & $0.007$ & $0.023$\\
 & $(3.59)$ & $(4.17)$ & $(2.81)$ & $(3.51)$ & $(2.20)$ & $(3.02)$ & $(1.98)$ & $(1.70)$ & $(0.24)$ & $(1.10)$\\
Option dummy & $0.057$ & $0.058$ & $0.080$ & $0.054$ & $0.022$ & $-0.054$ & $0.034$ & $0.057$ & $0.049$ & $-0.015$\\
 & $(0.70)$ & $(0.64)$ & $(0.93)$ & $(0.48)$ & $(0.20)$ & $(-0.51)$ & $(0.39)$ & $(0.63)$ & $(0.49)$ & $(-0.19)$\\
Amihud & $1.118$ & $0.352$ & $1.848$ & $2.474$ & $5.798^{**}$ & $5.028^{**}$ & $2.661$ & $5.542^{**}$ & $4.894^{**}$ & $3.491^{*}$\\
 & $(0.57)$ & $(0.17)$ & $(0.92)$ & $(0.89)$ & $(2.06)$ & $(2.03)$ & $(1.38)$ & $(2.55)$ & $(2.03)$ & $(1.96)$\\
May spread & $0.039$ & $0.025$ & $0.024$ & $-0.008$ & $-0.182^{***}$ & $-0.132^{***}$ & $-0.102^{***}$ & $-0.162^{***}$ & $-0.143^{***}$ & $-0.180^{***}$\\
 & $(1.46)$ & $(0.85)$ & $(0.80)$ & $(-0.21)$ & $(-4.65)$ & $(-3.78)$ & $(-3.44)$ & $(-4.89)$ & $(-3.93)$ & $(-6.65)$\\
Observations & 1,138 & 1,204 & 1,200 & 1,133 & 999 & 1,159 & 1,247 & 1,197 & 1,155 & 1,169\\
Adj.\ $R^2$ & .759 & .646 & .737 & .655 & .720 & .754 & .813 & .796 & .798 & .888\\
\bottomrule
\end{longtable}}
{\scriptsize\setlength{\tabcolsep}{2pt}
\begin{longtable}{@{}l*{10}{r}@{}}
\caption*{\textbf{Table 10, p.~30. Hypothetical arbitrage. Issuer fixed effects, clustered by issuer.}}\\
\toprule
 & Aug 13 & Sep 13 & Oct 13 & Nov 13 & Dec 13 & Jan 14 & Feb 14 & Mar 14 & Apr 14 & May 14\\
\midrule\endhead
ETFArb & $0.097^{*}$ & $0.148^{***}$ & $0.132^{**}$ & $0.144^{**}$ & $0.043$ & $0.091$ & $0.030$ & $-0.046$ & $-0.112^{*}$ & $-0.052$\\
 & $(1.72)$ & $(2.70)$ & $(2.43)$ & $(2.15)$ & $(0.57)$ & $(1.33)$ & $(0.52)$ & $(-0.76)$ & $(-1.69)$ & $(-0.87)$\\
Time to maturity & $-0.002$ & $0.001$ & $-0.002$ & $-0.002$ & $0.000$ & $0.007^{***}$ & $0.006^{***}$ & $0.011^{***}$ & $0.009^{***}$ & $0.013^{***}$\\
 & $(-1.49)$ & $(1.09)$ & $(-1.31)$ & $(-1.36)$ & $(0.22)$ & $(4.46)$ & $(5.12)$ & $(7.98)$ & $(6.16)$ & $(9.86)$\\
Coupon & $-0.016^{**}$ & $-0.030^{***}$ & $-0.036^{***}$ & $-0.049^{***}$ & $-0.049^{***}$ & $-0.063^{***}$ & $-0.044^{***}$ & $-0.043^{***}$ & $-0.048^{***}$ & $-0.049^{***}$\\
 & $(-2.24)$ & $(-4.12)$ & $(-5.27)$ & $(-5.62)$ & $(-4.91)$ & $(-7.10)$ & $(-6.26)$ & $(-5.34)$ & $(-5.45)$ & $(-6.25)$\\
log(size) & $0.000$ & $0.013$ & $-0.006$ & $-0.031$ & $0.004$ & $0.056^{*}$ & $0.036$ & $0.084^{***}$ & $0.057^{*}$ & $0.054^{*}$\\
 & $(0.00)$ & $(0.50)$ & $(-0.22)$ & $(-1.00)$ & $(0.11)$ & $(1.73)$ & $(1.37)$ & $(2.89)$ & $(1.81)$ & $(1.91)$\\
Option dummy & $-0.061$ & $-0.016$ & $-0.070$ & $-0.103$ & $-0.145$ & $-0.120$ & $-0.085$ & $-0.072$ & $-0.131$ & $-0.086$\\
 & $(-0.60)$ & $(-0.16)$ & $(-0.70)$ & $(-0.81)$ & $(-1.08)$ & $(-1.03)$ & $(-0.79)$ & $(-0.64)$ & $(-1.09)$ & $(-0.72)$\\
Amihud & $2.642$ & $3.584^{**}$ & $3.393^{**}$ & $4.297^{*}$ & $5.876^{**}$ & $7.505^{***}$ & $5.596^{***}$ & $7.061^{***}$ & $6.774^{***}$ & $4.862^{***}$\\
 & $(1.56)$ & $(2.07)$ & $(1.99)$ & $(1.86)$ & $(2.44)$ & $(3.59)$ & $(3.33)$ & $(3.72)$ & $(3.27)$ & $(2.65)$\\
May spread & $-0.013$ & $-0.075^{**}$ & $0.031$ & $-0.008$ & $-0.105^{***}$ & $-0.155^{***}$ & $-0.142^{***}$ & $-0.191^{***}$ & $-0.083^{**}$ & $-0.151^{***}$\\
 & $(-0.43)$ & $(-2.51)$ & $(1.08)$ & $(-0.23)$ & $(-2.68)$ & $(-4.44)$ & $(-4.77)$ & $(-5.66)$ & $(-2.23)$ & $(-4.70)$\\
Observations & 1,258 & 1,414 & 1,408 & 1,328 & 1,275 & 1,333 & 1,491 & 1,407 & 1,343 & 1391\\
Adj.\ $R^2$ & .856 & .848 & .904 & .933 & .936 & .907 & .951 & .938 & .881 & .909\\
\bottomrule
\end{longtable}}
{\scriptsize\setlength{\tabcolsep}{5pt}
\begin{longtable}{@{}llrrrrr@{}}
\caption*{\textbf{Table A1, p.~33. Monthly quintiles on net creation and redemption.}}\\
\toprule
Period & Quintile & Net (\$ mln) & Redemptions (\$ mln) & Avg daily dev. (\%) & Avg month-end dev. (\%) & Avg volatility\\
\midrule\endhead
May & $1$ & $-1{,}870.00$ & $-3{,}220.00$ & $0.0596$ & $-0.1751$ & $0.0022$\\
 & $2$ & $9.48$ & $0.00$ & $-0.0994$ & $-0.1093$ & $0.0032$\\
 & $3$ & $116.00$ & $0.00$ & $0.5340$ & $0.4654$ & $0.0023$\\
 & $4$ & $324.00$ & $0.00$ & $0.3081$ & $0.0852$ & $0.0013$\\
 & $5$ & $3{,}990.00$ & $-596.00$ & $0.1644$ & $0.0875$ & $0.0009$\\
June & $1$ & $-7{,}480.00$ & $-8{,}780.00$ & $-0.1925$ & $-0.0532$ & $0.0035$\\
 & $2$ & $-65.20$ & $-257.00$ & $-0.2749$ & $-0.1834$ & $0.0042$\\
 & $3$ & $0.00$ & $0.00$ & $-0.2172$ & $-0.4053$ & $0.0075$\\
 & $4$ & $136.00$ & $-6.46$ & $0.2914$ & $0.2084$ & $0.0028$\\
 & $5$ & $3{,}090.00$ & $-323.00$ & $0.2155$ & $0.2208$ & $0.0019$\\
July & $1$ & $-1{,}820.00$ & $-2{,}200.00$ & $-0.0737$ & $-0.0172$ & $0.0022$\\
 & $2$ & $-27.70$ & $-27.70$ & $-0.2878$ & $-0.4935$ & $0.0075$\\
 & $3$ & $7.99$ & $-22.60$ & $-0.0756$ & $-0.1721$ & $0.0041$\\
 & $4$ & $258.00$ & $-396.00$ & $0.2279$ & $0.2606$ & $0.0032$\\
 & $5$ & $5{,}010.00$ & $-814.00$ & $0.3027$ & $0.3498$ & $0.0016$\\
August & $1$ & $-4{,}350.00$ & $-4{,}990.00$ & $-0.1353$ & $-0.1078$ & $0.0017$\\
 & $2$ & $-18.60$ & $-35.00$ & $-0.6410$ & $-0.6610$ & $0.0039$\\
 & $3$ & $1.15$ & $0.00$ & $-0.2357$ & $-0.4126$ & $0.0058$\\
 & $4$ & $136.00$ & $-96.30$ & $0.3592$ & $0.3057$ & $0.0078$\\
 & $5$ & $2{,}630.00$ & $-300.00$ & $0.2616$ & $0.2590$ & $0.0009$\\
\bottomrule
\end{longtable}}
{\scriptsize\setlength{\tabcolsep}{2pt}
\begin{longtable}{@{}l*{10}{r}@{}}
\caption*{\textbf{Table B1, p.~34. Key rows only.}}\\
\toprule
 & Aug 13 & Sep 13 & Oct 13 & Nov 13 & Dec 13 & Jan 14 & Feb 14 & Mar 14 & Apr 14 & May 14\\
\midrule\endhead
\multicolumn{11}{@{}l}{\emph{A. Match on Roll}}\\
LQD bond & $0.138^{***}$ & $0.081^{**}$ & $0.104^{***}$ & $0.107^{***}$ & $0.071$ & $0.078^{**}$ & $0.010$ & $0.052$ & $0.038$ & $0.065$\\
 & $(3.54)$ & $(2.35)$ & $(3.58)$ & $(3.01)$ & $(1.31)$ & $(2.10)$ & $(0.30)$ & $(1.54)$ & $(0.75)$ & $(1.68)$\\
MF pressure & $0.025$ & $-0.019$ & $0.009$ & $0.039^{*}$ & $0.022$ & $0.051^{*}$ & $0.006$ & $0.017$ & $0.013$ & $0.015$\\
 & $(0.79)$ & $(-0.81)$ & $(0.41)$ & $(1.85)$ & $(0.55)$ & $(1.84)$ & $(0.24)$ & $(0.72)$ & $(0.38)$ & $(0.62)$\\
Roll & $12.158$ & $4.941$ & $-6.103$ & $10.960$ & $21.051$ & $28.277^{**}$ & $5.815$ & $16.781$ & $9.815$ & $8.274$\\
 & $(0.90)$ & $(0.45)$ & $(-0.67)$ & $(1.05)$ & $(1.66)$ & $(2.67)$ & $(0.68)$ & $(1.58)$ & $(0.78)$ & $(0.86)$\\
Observations & 151 & 161 & 160 & 152 & 148 & 158 & 168 & 167 & 161 & 161\\
Adj.\ $R^2$ & .901 & .915 & .946 & .936 & .899 & .915 & .836 & .874 & .831 & .896\\\midrule
\multicolumn{11}{@{}l}{\emph{B. Match on IRC}}\\
LQD bond & $0.094$ & $0.104^{***}$ & $0.083^{**}$ & $0.099^{***}$ & $0.061$ & $0.096^{***}$ & $0.010$ & $0.029$ & $-0.036$ & $0.003$\\
 & $(1.62)$ & $(3.88)$ & $(2.58)$ & $(2.94)$ & $(1.29)$ & $(2.83)$ & $(0.30)$ & $(0.77)$ & $(-0.85)$ & $(0.07)$\\
MF pressure & $0.049$ & $0.062^{**}$ & $0.021$ & $0.056^{**}$ & $0.015$ & $0.052^{*}$ & $-0.004$ & $-0.022$ & $-0.024$ & $0.031$\\
 & $(1.32)$ & $(2.76)$ & $(0.76)$ & $(2.09)$ & $(0.36)$ & $(1.94)$ & $(-0.15)$ & $(-0.80)$ & $(-0.78)$ & $(0.93)$\\
IRC & $20.545$ & $36.919^{**}$ & $22.429$ & $40.956$ & $41.939$ & $42.573^{*}$ & $5.987$ & $31.410$ & $16.383$ & $24.938$\\
 & $(0.48)$ & $(2.11)$ & $(0.97)$ & $(1.70)$ & $(1.47)$ & $(1.93)$ & $(0.32)$ & $(1.16)$ & $(0.67)$ & $(1.13)$\\
Observations & 129 & 141 & 140 & 135 & 133 & 137 & 148 & 145 & 145 & 139\\
Adj.\ $R^2$ & .946 & .975 & .960 & .969 & .960 & .967 & .923 & .916 & .916 & .944\\
\bottomrule
\end{longtable}}
{\scriptsize\setlength{\tabcolsep}{2pt}
\begin{longtable}{@{}l*{10}{r}@{}}
\caption*{\textbf{Table B2, p.~35. Match on time to maturity. Key rows only.}}\\
\toprule
 & Aug 13 & Sep 13 & Oct 13 & Nov 13 & Dec 13 & Jan 14 & Feb 14 & Mar 14 & Apr 14 & May 14\\
\midrule\endhead
\multicolumn{11}{@{}l}{\emph{A. Amihud}}\\
LQD bond & $0.168^{***}$ & $0.095^{**}$ & $0.090^{**}$ & $0.109^{**}$ & $0.062$ & $0.056$ & $-0.018$ & $0.038$ & $0.072$ & $0.065^{*}$\\
 & $(4.07)$ & $(2.31)$ & $(2.33)$ & $(2.69)$ & $(1.04)$ & $(1.66)$ & $(-0.41)$ & $(1.15)$ & $(1.65)$ & $(1.75)$\\
MF pressure & $-0.046$ & $-0.039$ & $-0.054^{*}$ & $0.011$ & $-0.084^{*}$ & $0.010$ & $-0.040$ & $-0.012$ & $0.022$ & $0.008$\\
 & $(-1.08)$ & $(-1.31)$ & $(-1.77)$ & $(0.38)$ & $(-1.98)$ & $(0.38)$ & $(-1.30)$ & $(-0.49)$ & $(0.64)$ & $(0.27)$\\
Amihud & $7.883^{*}$ & $2.382$ & $1.429$ & $9.075^{**}$ & $8.950^{*}$ & $11.378^{***}$ & $5.319$ & $11.565^{***}$ & $9.954^{***}$ & $8.994^{***}$\\
 & $(2.03)$ & $(0.60)$ & $(0.37)$ & $(2.14)$ & $(1.93)$ & $(3.71)$ & $(1.50)$ & $(3.27)$ & $(2.80)$ & $(2.97)$\\
Observations & 130 & 153 & 152 & 147 & 145 & 149 & 165 & 164 & 157 & 152\\
Adj.\ $R^2$ & .895 & .846 & .852 & .875 & .851 & .901 & .800 & .898 & .877 & .909\\\midrule
\multicolumn{11}{@{}l}{\emph{B. Roll (October = November, March = April, as printed)}}\\
LQD bond & $0.183^{***}$ & $0.111^{**}$ & $0.085^{*}$ & $0.085^{*}$ & $0.087$ & $0.070$ & $0.002$ & $0.044$ & $0.044$ & $0.096^{*}$\\
 & $(4.03)$ & $(2.33)$ & $(1.85)$ & $(1.85)$ & $(1.13)$ & $(1.50)$ & $(0.03)$ & $(0.99)$ & $(0.99)$ & $(1.98)$\\
MF pressure & $-0.055$ & $-0.073^{*}$ & $-0.050$ & $-0.050$ & $-0.145^{**}$ & $0.024$ & $-0.078^{*}$ & $-0.013$ & $-0.013$ & $-0.021$\\
 & $(-1.25)$ & $(-1.79)$ & $(-1.31)$ & $(-1.31)$ & $(-2.15)$ & $(0.60)$ & $(-1.82)$ & $(-0.36)$ & $(-0.36)$ & $(-0.46)$\\
Roll & $22.497$ & $2.284$ & $-18.945$ & $-18.945$ & $3.856$ & $28.158^{*}$ & $11.511$ & $19.915$ & $19.915$ & $18.056$\\
 & $(1.30)$ & $(0.12)$ & $(-1.05)$ & $(-1.05)$ & $(0.20)$ & $(1.88)$ & $(0.86)$ & $(1.23)$ & $(1.23)$ & $(1.36)$\\
Observations & 130 & 153 & 152 & 147 & 145 & 149 & 165 & 164 & 157 & 152\\
Adj.\ $R^2$ & .886 & .850 & .856 & .856 & .839 & .877 & .782 & .876 & .876 & .882\\\midrule
\multicolumn{11}{@{}l}{\emph{C. IRC}}\\
LQD bond & $0.202^{***}$ & $0.126^{**}$ & $0.103^{**}$ & $0.111^{***}$ & $0.105$ & $0.062$ & $0.018$ & $0.061$ & $0.080$ & $0.072$\\
 & $(4.96)$ & $(2.67)$ & $(2.36)$ & $(2.85)$ & $(1.56)$ & $(1.45)$ & $(0.37)$ & $(1.45)$ & $(1.53)$ & $(1.64)$\\
MF pressure & $-0.048$ & $-0.071^{*}$ & $-0.055$ & $0.003$ & $-0.122^{*}$ & $0.030$ & $-0.075^{*}$ & $-0.005$ & $0.038$ & $-0.010$\\
 & $(-1.22)$ & $(-1.77)$ & $(-1.47)$ & $(0.09)$ & $(-1.99)$ & $(0.81)$ & $(-1.83)$ & $(-0.15)$ & $(0.91)$ & $(-0.25)$\\
IRC & $73.320^{***}$ & $37.958$ & $40.289$ & $90.417^{***}$ & $127.291^{**}$ & $78.948^{**}$ & $62.310^{*}$ & $69.453^{**}$ & $162.955^{***}$ & $83.189^{***}$\\
 & $(2.93)$ & $(1.27)$ & $(1.13)$ & $(3.33)$ & $(2.43)$ & $(2.65)$ & $(1.85)$ & $(2.64)$ & $(4.53)$ & $(3.01)$\\
Observations & 130 & 153 & 152 & 147 & 145 & 149 & 165 & 164 & 157 & 152\\
Adj.\ $R^2$ & .908 & .857 & .857 & .917 & .870 & .888 & .796 & .891 & .904 & .903\\
\bottomrule
\end{longtable}}
{\scriptsize\setlength{\tabcolsep}{2pt}
\begin{longtable}{@{}l*{10}{r}@{}}
\caption*{\textbf{Table B3, p.~36. Key rows only.}}\\
\toprule
 & Aug 13 & Sep 13 & Oct 13 & Nov 13 & Dec 13 & Jan 14 & Feb 14 & Mar 14 & Apr 14 & May 14\\
\midrule\endhead
\multicolumn{11}{@{}l}{\emph{A. Roll}}\\
ETF pressure & $0.012$ & $0.040^{**}$ & $0.041^{***}$ & $0.044^{**}$ & $0.046^{**}$ & $0.038^{**}$ & $0.014$ & $0.020$ & $-0.007$ & $0.002$\\
 & $(0.78)$ & $(2.54)$ & $(2.69)$ & $(2.29)$ & $(2.29)$ & $(2.04)$ & $(0.91)$ & $(1.20)$ & $(-0.41)$ & $(0.17)$\\
MF pressure & $-0.018^{***}$ & $-0.011^{*}$ & $-0.003$ & $0.005$ & $-0.006$ & $-0.003$ & $0.010$ & $0.008$ & $0.007$ & $0.006$\\
 & $(-2.71)$ & $(-1.69)$ & $(-0.39)$ & $(0.56)$ & $(-0.66)$ & $(-0.32)$ & $(1.57)$ & $(1.16)$ & $(0.94)$ & $(0.87)$\\
Roll & $4.436$ & $16.253^{**}$ & $9.149$ & $16.350^{*}$ & $21.685^{**}$ & $11.446$ & $2.594$ & $10.474$ & $-1.503$ & $-4.159$\\
 & $(0.69)$ & $(2.28)$ & $(1.25)$ & $(1.87)$ & $(2.11)$ & $(1.37)$ & $(0.43)$ & $(1.40)$ & $(-0.20)$ & $(-0.64)$\\
Observations & 1,138 & 1,204 & 1,200 & 1,133 & 999 & 1,159 & 1,247 & 1,197 & 1,155 & 1,169\\
Adj.\ $R^2$ & .759 & .648 & .737 & .656 & .720 & .753 & .813 & .795 & .797 & .887\\\midrule
\multicolumn{11}{@{}l}{\emph{B. IRC}}\\
ETF pressure & $0.012$ & $0.034^{**}$ & $0.036^{**}$ & $0.045^{**}$ & $0.043^{**}$ & $0.034^{*}$ & $0.038^{**}$ & $0.023$ & $-0.006$ & $0.000$\\
 & $(0.83)$ & $(2.29)$ & $(2.55)$ & $(2.49)$ & $(2.44)$ & $(1.94)$ & $(2.39)$ & $(1.32)$ & $(-0.34)$ & $(0.01)$\\
MF pressure & $-0.019^{***}$ & $-0.011$ & $-0.002$ & $0.004$ & $-0.002$ & $-0.001$ & $-0.006$ & $-0.011$ & $0.009$ & $0.006$\\
 & $(-2.85)$ & $(-1.63)$ & $(-0.27)$ & $(0.55)$ & $(-0.19)$ & $(-0.08)$ & $(-0.79)$ & $(-1.34)$ & $(1.12)$ & $(0.92)$\\
IRC & $21.684^{*}$ & $4.071$ & $11.947$ & $31.015^{**}$ & $23.908$ & $28.474^{**}$ & $26.148^{**}$ & $29.767^{**}$ & $7.356$ & $0.833$\\
 & $(1.86)$ & $(0.35)$ & $(1.01)$ & $(1.96)$ & $(1.53)$ & $(2.03)$ & $(2.14)$ & $(2.17)$ & $(0.51)$ & $(0.08)$\\
Observations & 1,137 & 1,261 & 1,254 & 1,190 & 1,157 & 1,205 & 1,299 & 1,258 & 1,199 & 1,226\\
Adj.\ $R^2$ & .760 & .761 & .825 & .778 & .839 & .838 & .850 & .853 & .865 & .925\\
\bottomrule
\end{longtable}}
{\scriptsize\setlength{\tabcolsep}{2pt}
\begin{longtable}{@{}l*{10}{r}@{}}
\caption*{\textbf{Table B4, p.~37. Key rows only.}}\\
\toprule
 & Aug 13 & Sep 13 & Oct 13 & Nov 13 & Dec 13 & Jan 14 & Feb 14 & Mar 14 & Apr 14 & May 14\\
\midrule\endhead
\multicolumn{11}{@{}l}{\emph{A. Roll}}\\
ETFArb & $0.096^{*}$ & $0.138^{**}$ & $0.122^{**}$ & $0.126^{*}$ & $0.031$ & $0.076$ & $0.011$ & $-0.020$ & $-0.085$ & $-0.058$\\
 & $(1.70)$ & $(2.56)$ & $(2.24)$ & $(1.86)$ & $(0.39)$ & $(1.09)$ & $(0.19)$ & $(-0.34)$ & $(-1.32)$ & $(-0.94)$\\
Roll & $3.649$ & $12.104^{**}$ & $4.566$ & $8.574$ & $15.247^{*}$ & $17.271^{**}$ & $7.210$ & $10.016^{*}$ & $7.128$ & $5.320$\\
 & $(0.67)$ & $(2.50)$ & $(0.78)$ & $(1.33)$ & $(1.88)$ & $(2.41)$ & $(1.48)$ & $(1.76)$ & $(1.23)$ & $(0.87)$\\
Observations & 1,258 & 1,414 & 1,408 & 1,328 & 1,275 & 1,333 & 1,491 & 1,407 & 1,343 & 1,391\\
Adj.\ $R^2$ & .856 & .860 & .906 & .933 & .936 & .906 & .951 & .941 & .891 & .908\\\midrule
\multicolumn{11}{@{}l}{\emph{B. IRC}}\\
ETFArb & $0.100^{*}$ & $0.133^{**}$ & $0.123^{**}$ & $0.123^{*}$ & $0.024$ & $0.078$ & $0.020$ & $-0.019$ & $-0.084$ & $-0.056$\\
 & $(1.77)$ & $(2.48)$ & $(2.26)$ & $(1.80)$ & $(0.31)$ & $(1.12)$ & $(0.35)$ & $(-0.32)$ & $(-1.30)$ & $(-0.91)$\\
IRC & $20.791^{**}$ & $13.899$ & $11.651$ & $3.490$ & $23.011$ & $34.178^{***}$ & $26.012^{***}$ & $27.630^{***}$ & $13.675$ & $12.185$\\
 & $(1.98)$ & $(1.48)$ & $(1.21)$ & $(0.29)$ & $(1.61)$ & $(2.79)$ & $(2.72)$ & $(2.70)$ & $(1.19)$ & $(1.10)$\\
Observations & 1,258 & 1,414 & 1,408 & 1,328 & 1,275 & 1,333 & 1,491 & 1,407 & 1,343 & 1,391\\
Adj.\ $R^2$ & .856 & .859 & .906 & .933 & .936 & .906 & .951 & .942 & .891 & .908\\
\bottomrule
\end{longtable}}

\subsection{Figure values, read by eye}
\IN{the paper prints no figure values, so these are my approximate readings. Use them only as rough checks}
\begin{itemize}
\item Figure 2 panel C, Low Deviation redemptions, \$ billions: May 2.6, June 6.2, July 2.2, August 3.4.
\item Figure 2 panel C, Discount redemptions, \$ billions: May 0.35, June 1.7, July 1.1, August 1.0.
\item Figure 2 panel C, Premium redemptions, \$ billions: May 0.9, June 1.5, July 0.2, August 1.0.
\item Figure 2 panel B, total assets: Low Deviation near \$110 billion, Discount near \$20 to \$25 billion, Premium near \$30 to \$34 billion.
\item Figure 3 panel A, LQD minus non-LQD spread: about $-0.025$ in January 2013, a peak near $0.103$ in August 2013, and near zero by July 2014.
\end{itemize}

\subsection{Quoted figures in the text, checked against the tables}
\begin{longtable}{@{}P{5.3cm}P{1.5cm}P{6.2cm}P{3.5cm}@{}}
\toprule
Text claim & Page & Arithmetic & Verdict\\
\midrule\endfirsthead
\toprule Text claim & Page & Arithmetic & Verdict\\ \midrule\endhead
ETF flow volatility 1.9 times mutual funds & 12 & $7.13/3.75=1.90$. $7.13/3.69=1.93$. & Matches\\
1.4 points above, a 37\% increase & 2 & $1.4/3.75=37.3\%$. Table 2 column 2 is 1.454. & Matches when rounded\\
Minimum 40\% (18.4\%) increase & 13 & $1.454/3.75=38.8\%$. $0.9/4.88=18.4\%$. & 40\% does not match 37\% on p.~2\\
Interaction gives 0.725\% and 0.790\% & 14 & $0.484+0.241=0.725$. $0.524+0.266=0.790$. & Matches\\
A 37.5\% and 41.1\% increase over mean flow & 14 & $0.725/1.92=37.8\%$. $0.790/1.92=41.1\%$. & First is off by 0.3\\
ETF outflow multiples of 9.3, 9.1, 7.6 & 19 & $7.82/0.84$, $7.66/0.84$, $6.37/0.84$. & Matches\\
Volume up 34\% and 15\% & 21 & $181{,}785/135{,}797$, $673{,}265/587{,}130$. & Matches\\
11.3 and 8.5 bps, as 10.1\% and 7.6\% of the spread & 26 & Both imply a pre-episode mean LQD spread near 1.12\%. & Consistent\\
4.0 bps in the month after & 5 & Table 9 September is 0.039, so 3.9 bps. & Text rounds up\\
14.8 bps & 5, 28 & Table 10 September is 0.148. & Matches\\
20.35 bps at peak & 5, 28 & $4.5\times 4.52=20.34$. & Matches within rounding\\
Average monthly ETF return is 22 bps & 28 & Table 1 mean is 0.20\%. & Does not match\\
Median ETF 3\%, active 15\%, index 43\% liquid & 3 & $0.91+2.13$, $3.77+11.24$, $1.33+42.06$. & Matches\\
Index median 42.6\% Treasuries, 26.23\% corporate & 16 & Table 1 has 42.06 and 26.38. & Does not match\\
Table 4 p-values & 18 & $0.946$, $t=15.13$ gives $|1-\beta|/\mathrm{se}=0.86$, $p\approx.39$. $0.872$, $t=15.75$ gives $2.31$, $p\approx.02$. & Matches\\
\bottomrule
\end{longtable}

% ------------------------------------------------------------------
\section{Label register}\label{sec:register}

Labels sit next to each item throughout this document. This section collects the ones that matter. The label counts in this file are: 180 STATED, 43 INFERRED, 67 MISSING and 9 MINE. These counts include the four examples in the reading guide.

\subsection{Gaps ranked by how much they could move Table 8, 9 or 10}\label{sec:gaps}
\begin{longtable}{@{}P{0.5cm}P{4.8cm}P{5.4cm}P{5.8cm}@{}}
\caption*{\textbf{Ranked gaps and conflicts}}\\
\toprule
\# & Gap & What the paper says & Proposed default\\
\midrule\endfirsthead
\toprule \# & Gap & What the paper says & Proposed default\\ \midrule\endhead
1 & Control bond rule & Closest in Amihud (p.~25, Table 7 note). Largest below the threshold (Table 8 note). These conflict. & \MI{closest May Amihud, because the body text and Table B1 match on liquidity. Also run the largest-below rule, and report which is closer to the targets}\\
2 & Which bonds count in Table 9 & ``In baskets exchanged with APs'' (p.~26). Table 9 has fewer bonds than Table 10. & \MI{require at least one ETF holding change}\\
3 & Spread construction & Volume-weighted yield over the matched Treasury rate. No source for yields or Treasury rates. & \MI{trade-level yield minus H.15 interpolated, then par-weighted mean}\\
4 & Winsorising scope & 1\% and 99\%, no scope given. & \MI{by month over all filtered bonds, before differencing}\\
5 & May-level Amihud window & ``May levels''. Onset falls inside May. & \MI{1 to 21 May}\\
6 & Whether treatment bonds must be held by LQD & Eligibility language, with no holding check. & \MI{eligibility first, holding as a variant}\\
7 & Index issuer rules & Cited by a link only. & \MI{size and maturity rules, then add the 2013 issuer rules if found}\\
8 & Amount-change window & No dates given. & \MI{22 May 2013 to 31 May 2014}\\
9 & Rating aggregation & ``Most recent'' from three agencies. & \MI{mean}\\
10 & Quarterly flow, new positions (Table 4) & Not defined. & \MI{Eq.~(1) between reports, and drop new positions}\\
11 & Table 7 turnover & Not defined. & \MI{monthly par turnover}\\
12 & Summary statistics of ETFArb & None reported. & \MI{compute and publish them}\\
\bottomrule
\end{longtable}

\subsection{Internal conflicts in the paper}
These are errors or tensions inside the paper. I checked each against the page image. None changes the headline result by itself. Items 1 and 2 change what you build. Item 3 means you cannot use those columns as targets.
\begin{enumerate}
\item \textbf{Control rule.} See the first row of the ranked table in Section~\ref{sec:gaps}. \ST{p.~25, p.~26, p.~27}
\item \textbf{Size control in Table 8.} The p.~26 text excludes it and the table shows no size row. The Table 8 note still lists the log of amount outstanding. I follow the text.
\item \textbf{Table B2 panel B has duplicated columns.} October equals November, and March equals April, in every coefficient and t-statistic. Bond counts differ (152 and 147, 164 and 157). The authors' supplementary file has the same numbers, so this is not a typesetting error. Treat those four columns as unreliable.
\item \textbf{Trading days.} ``Last three full trading days'' (p.~8, p.~25) against ``last 3 business days'' (Tables 8--10 notes).
\item \textbf{No negative reversal claim.} P.~31 says the pattern of negative significant coefficients after September appears in no regression. Table 10 April 2014 is $-0.112$, $t=-1.69$, significant at 10\%. It is a single month at the weakest star level.
\item \textbf{Table 7 signs.} MF pressure, Amihud and turnover have negative differences with positive t-values. Time to maturity and coupon both show $t=1.14$, with p-values .269 and .254.
\item \textbf{Amihud scale.} Table 1 mean is 0.73, and Table 7 means are 0.011 and 0.009.
\item \textbf{Figure 2 dates.} P.~21 says January 2013 to June 2014. The caption and the axis say January 2012 to December 2014.
\item \textbf{Figure 1 axes.} The caption swaps the axes of panel B. The figure's axis titles match the p.~6 text.
\item \textbf{Redemptions share.} P.~21 says Low Deviation ETFs account for ``nearly all'' redemptions. Reading the bars of panel C by eye gives about two-thirds. P.~4 says ``most''. \IN{from my reading of the bars, so this is approximate}
\item \textbf{Table 6 volume.} The row says shares and the note says dollars.
\item \textbf{Notes.} The Table 9 note defines the option dummy as ``equal to zero if a bond has option features and zero otherwise''. The Table 10 note says it repeats Table 7, and it should say Table 8 or 9.
\item \textbf{LQD history.} Fn.~11 says 18-year history. Fn.~26 says 126-month history through June 2020.
\item \textbf{Bond counts across proxies.} Table B3 panel A has the same counts as Table 9. Table B4 panels A and B have the same counts as Table 10. The liquidity proxies should have different missing patterns. Table B3 panel B counts do differ. \IN{the identical counts suggest a shared sample or missing values filled in, which the paper does not explain}
\item \textbf{December 2013 count.} Table 9 has 999 bonds in December, against 1,133 to 1,247 in other months. Table B3 panel B has 1,157 in that month. The paper gives no reason.
\end{enumerate}

\subsection{Errors in the earlier notes}
I checked \texttt{E\_local\_documents.md}, Section 4.4 (lines 117--136). \textbf{I found no wrong number. The notes are incomplete in five places.}
\begin{enumerate}
\item \textbf{``8.5 to 11.3 bps''.} The note gives 11.3 in August and 8.5 in September, which is what the paper says on p.~26. But October (11.2) and November (10.2) are also significant. So 8.5 is the low of a four-month range, not the end of a decay. The path does not fall steadily. It is also not a confidence interval.
\item \textbf{``Episode runs 21 May to 18 September''.} Only the LQD share statistic uses those dates. The paper's onset is 22 May, and the pressure measures stop at the end of August.
\item \textbf{``4.0 bps a month after the episode''.} The text on p.~5 says this. Table 9 gives 0.039, so 3.9 bps. Use 3.9.
\item \textbf{``1.4 percentage points, 37\%''.} This follows p.~2. Table 2 column 2 is 1.454, and p.~13 says a minimum of 40\%. The p.~2 phrase ``lagged fund characteristics'' also does not fit column 2, whose controls are sample averages. The lagged panel, column 4, gives 2.155.
\item \textbf{``One for one''.} This holds only for outflows and only with controls: 0.946, $p=.39$. Without controls the slope is 0.872, and the paper's table rejects one at $p=.02$. Inflows are 0.411. The note's data list also omits Morningstar, Bloomberg and Mergent FISD.
\end{enumerate}

% ------------------------------------------------------------------
\section{Threats to validity}\label{sec:threats}

\subsection{As the paper handles them}
\begin{longtable}{@{}P{5.2cm}P{8.8cm}P{3cm}@{}}
\toprule
Threat & Paper's response & Label\\
\midrule\endhead
Investors sell the funds most exposed to rates. Managers sell the weakest bonds. & Issuer fixed effects compare bonds with one credit risk. & \ST{p.~4, p.~23}\\
Selection into ETFs. & Only bonds in ETF portfolios enter. & \ST{p.~4, p.~24}\\
Term structure differences. & Spreads against maturity-matched Treasury rates. & \ST{p.~4, p.~25}\\
Selling due to the 3-year index rule. & At least 4 years to maturity in setting 1. & \ST{p.~4, p.~25}\\
Selling due to the 1-year index rule. & Drop bonds under 1 year and 4 months. & \ST{p.~11}\\
Forced selling on a downgrade. & Drop bonds crossing investment grade, by any one agency. & \ST{p.~11, fn.~15}\\
Treatment and control differ. & Table 7 shows no significant gap in rating, maturity, coupon, MF pressure, Amihud or turnover. & \ST{p.~26}\\
Mutual fund selling drives the result. & MF pressure control. It is insignificant in Table 8 and negative in early Table 9 months. & \ST{p.~26, p.~28}\\
Manager discretion in baskets. & Hypothetical measure with pre-episode weights. & \ST{p.~28}\\
Mechanical link in scaling. & Scale by amount outstanding, not dollar volume, citing Wardlaw (2020). & \ST{p.~12}\\
Choice of liquidity proxy or matching variable. & Tables B1--B4. & \ST{p.~26, p.~28}\\
Price discovery. & Coefficient path and Figure 3 panel B. & \ST{p.~5, p.~31}\\
Timing of AP sales unobserved. & Acknowledged. Dealer inventory half-life of 4 to 6 weeks or over a month supports a slow effect. & \ST{fn.~7, p.~28}\\
One episode. & Discussed through a contrast with COVID-19, where APs let discounts persist and the Fed intervened. & \ST{p.~6, p.~32}\\
\bottomrule
\end{longtable}

\subsection{My own concerns}
All items below are \MY{my analysis, using the paper's own numbers. They are not claims by the authors}
\begin{enumerate}
\item \textbf{Table 8 may partly measure bond size.} This is the most important concern. Treatment bonds average \$1,122 million and controls \$416 million. So this is not a narrow comparison at the threshold. The log gap is 0.99. The Table 9 slopes on $\log$ size in August to November are 0.086, 0.103, 0.069 and 0.105. Those slopes alone imply 7 to 10 bps across that gap, close to the LQD coefficients of 8.5 to 11.3 bps. Table 8 omits size. But the evidence is mixed: the Table 10 size slopes are near zero in those months. And in Table 9 the size slope may itself capture LQD eligibility.
\item \textbf{The measured pressure gap does not explain the LQD effect.} Table 7 shows a pressure gap of 0.181. At the Table 9 peak slope of 0.045, that predicts 0.8 bps. Table 8 shows 8.5 to 11.3 bps, more than ten times larger. Either the pressure measure has large errors, or the LQD dummy captures something else, such as size or LQD-specific trading.
\item \textbf{No test of pre-trends.} Panel A of Figure 3 shows the LQD gap at about $-2.5$ bps in January and February 2013, rising to about $+2$ to $+3$ bps in March and April, before the onset. These are my readings of the chart. The authors attribute early moves to market jitters (p.~23). A placebo regression with a pre-May baseline would test this. The May baseline may also already contain some of the effect, which would understate it.
\item \textbf{Repeated cross-sections.} A bond enters a month only if it trades in the last three full days. So the sample changes month to month, and trading itself may respond to pressure. A fading coefficient could partly reflect a change in sample. A balanced panel would test this.
\item \textbf{Many tests.} Tables 8--10 hold 30 monthly tests. No correction for multiple testing appears. Most significant t-statistics lie between 2.05 and 2.70. The months overlap in sample and baseline, so the tests are not independent.
\item \textbf{Few clusters in Table 8.} I infer about 75 to 85 issuers. With issuer effects absorbing about half of the 151 to 171 observations, little variation remains. Adjusted R-squared of .837 to .904 may mostly reflect the fixed effects. Wild cluster bootstrap inference would be a useful check.
\item \textbf{Regression to the mean.} The May spread sits on both sides of the equation. Its coefficient is often large and negative. If May spreads are noisy, the control absorbs noise and may bias other slopes.
\item \textbf{Yield spread is not bond-specific in rate risk.} Treasury yields rose sharply. The maturity match leaves convexity, coupon and call effects in the spread. Table 8 treated and control bonds average 15 to 16 years, so these effects can be large. An option-adjusted measure would be cleaner. This also weakens the claim that the difference is ``in excess of'' term structure changes.
\item \textbf{The channel is inferred, not observed.} The paper does not show that APs sold these bonds. Bond-level customer selling in TRACE, by month and dealer direction, could test it directly.
\item \textbf{Holdings timing.} Pressure spans report dates that can sit well outside May to August. Index rebalancing and ordinary trading enter $\Delta H$.
\item \textbf{The August regression is not out of sample.} The ETF measures run to 31 August. The August spread is measured in the same days.
\item \textbf{Reversal is confounded with the market.} Spreads fall for all bonds after September (Figure 3 panel B). Reversal coincides with a market-wide rally and the 18 September Fed decision. This is the authors' setting, so it cannot be separated within one episode.
\item \textbf{External validity.} This is one rate-driven episode in a young market. ETF ownership averaged about 2\% per bond (Table 1). ETF assets have since more than doubled (p.~6, p.~32). AP behaviour differed in March 2020 (p.~6, p.~32). Rules on ETF baskets may also have changed since 2013, which I did not check. A credit-driven episode, or a calm period, may behave differently.
\end{enumerate}

% ------------------------------------------------------------------
\section{Adaptation to TRACE positioning, ETF holdings or RFQ flow}\label{sec:adapt}

This whole section is \MY{my view, not the paper's, and it is not a study design}

\textbf{What would carry over.}
\begin{itemize}
\item \textbf{Same-issuer comparison.} Issuer fixed effects remove issuer credit news from any bond-level flow measure. Any TRACE or RFQ pressure signal needs this control.
\item \textbf{Actual against hypothetical exposure.} The trade-based and weight-based pair is a useful template. Weights times flows avoid manager choice. They work wherever holdings and shares outstanding exist.
\item \textbf{Creation and redemption dollars as flow.} Eq.~(4) needs only NAV and shares. It extends to any ETF with daily data.
\item \textbf{The reversal test.} A pressure signal should predict a move and then its reversal over months. So a 10-day label sees only the first part.
\item \textbf{Scaling by amount outstanding.} This avoids the mechanical link with volume.
\item \textbf{Last-three-days spread.} It is a simple way to get month-end spreads for illiquid bonds.
\item \textbf{The deviation sort.} The result that the largest ETFs carry the redemptions suggests weighting flow by the ETF's arbitrage activity.
\end{itemize}

\textbf{What would not carry over.}
\begin{itemize}
\item \textbf{The event design.} The design needs one large shock. A continuous signal needs a panel with time effects and a different inference approach.
\item \textbf{The sample size with 15 issuers.} Issuer effects with 15 issuers leave too few clusters for the paper's inference, and no issuer universe to generalise from.
\item \textbf{The LQD threshold.} The \$750 million rule belongs to the 2013 index. Any current use needs today's rule book. Most of the 15 issuers may have few bonds below it.
\item \textbf{Quarterly holdings.} These are too slow for RFQ-scale signals. N-PORT and daily sponsor files change what is possible, but not for 2013.
\item \textbf{Causal reading of RFQ flow.} The paper's flows come from a fund's redemptions, which are plausibly outside the bond's own news. RFQ flow is chosen bond by bond, so it may carry bond-specific information. The hypothetical measure has no direct RFQ analogue.
\item \textbf{Dealer direction.} The paper uses no buy or sell side from TRACE and no dealer data. So it gives no guidance on signing flow.
\item \textbf{Magnitudes.} The basis point effects belong to 2013 ownership levels and that one episode.
\end{itemize}

% ------------------------------------------------------------------
\section{Acceptance checks, cheapest first}\label{sec:accept}

The tolerances are \MI{within 1 standard error for key coefficients, within 5\% for bond counts and within 10\% for fund counts. Same star level or one level away. The paper sets no tolerance}

\begin{enumerate}
\item \textbf{Code test with synthetic data.} Plant a known $\beta_1$ in a fake panel with issuer effects. Recover it. Check that flipping the sign flips $\beta_1$. Cost: no data needed.
\item \textbf{Arithmetic of the targets.} Reproduce the table of quoted figures in Section~\ref{sec:targets}. Cost: no data needed.
\item \textbf{LQD share count.} Shares outstanding fall 22.9\% from 21 May to 18 September 2013. Average daily deviation is 0.1 bps, and the worst end-of-day discount is 87 bps (p.~23). Cost: public sponsor data.
\item \textbf{Table 5.} Reproduce the 10th percentile flows and the multiples 9.3, 9.1 and 7.6 (p.~19). Cost: CRSP monthly only.
\item \textbf{Table 1 panel A.} Match means and medians. Check 7.13 against 3.75 for flow volatility. Cost: CRSP and Morningstar.
\item \textbf{Tables 2 and 3.} Match the ETF coefficient in columns 1 and 2 of Table 2, then the interaction in Table 3 columns 1 and 4. Cost: CRSP and Morningstar.
\item \textbf{Redemptions tie out.} Table A1 redemptions by month sum to about \$3.8, \$9.4, \$3.5 and \$5.4 billion. Figure 2 panel C by eye gives the same totals. \IN{I summed the Table A1 rows and read the bars} Cost: daily ETF data.
\item \textbf{Table 6 and Figure 2.} Low Deviation means near 3,795 and Discount near 905 in pre-Tantrum size. Cost: daily ETF data.
\item \textbf{Table 4, columns without controls.} ETF outflow 0.872 and $p=.02$. Inflow 0.362. Cost: holdings.
\item \textbf{Bond universe counts.} 10,131 bonds, 67\%, 96\% and 48\%, 1,491 issuers and 68. Cost: holdings and FISD.
\item \textbf{Table 4 with controls.} ETF outflow 0.946 and $p=.39$. Cost: adds full TRACE.
\item \textbf{Scale check.} May Amihud means near 0.011 and 0.009 in the Table 7 sample, and Table 1B Amihud near 0.73. If you get these, your scale matches the paper's two scales. Cost: TRACE 2013.
\item \textbf{Table 10, then Table 9.} Run these before Table 8 because they have no matching step. Targets: $N$, then the sign pattern. Table 10 is positive and significant August to November. Table 9 is significant September to January. Cost: all of the event data.
\item \textbf{Table 7, then Table 8.} Build both control rules. The right rule should give 151 to 171 bonds, the Table 7 balance and 0.113, 0.085, 0.112, 0.102 for August to November. Cost: event data plus the index rules.
\item \textbf{Appendix B.} Compare with Tables B1, B3 and B4. Ignore the four duplicated columns in B2 panel B.
\item \textbf{Falsification beyond the paper.} \MY{run the same regressions with a February 2013 baseline and April 2013 as the outcome. The ETF proxies should have zero slope. A non-zero slope points to the pre-trend concern in Section~\ref{sec:threats}}
\end{enumerate}

\textbf{What failure means.} A gap in checks 1 to 12 is a build error. Fix it before going further. A gap in 13 or 14 may be a build error, or the result of a Section~\ref{sec:register} choice. Vary those choices one at a time, in ranked order, and log which ones close the gap.

\end{document}
